\chapter{{Postnominal} {and} {pronominal} {\textit{cualquiera} }{in} {Argentinian} {Spanish}}\label{sec:4}

This chapter analyses \textit{cualquiera} in Argentinian Spanish and accounts for the differences from its counterpart in European Spanish.\footnote{{Aspects of this topic were presented in my talk “}{\textit{Cualquiera}} {in Argentinian Spanish”, at II Encuentro de Dialectos del Español (Ciudad Real, Universidad de Castilla-La Mancha, 17–18 May 2018). I’m very thankful to the audience of this workshop, especially to Carlos Muñoz.}} I will suggest that \textit{cualquiera} embedded under certain transitive verbs such as ArgSp \textit{mandar/decir cualquiera} ‘talk nonsense’ has its origin in \textit{mandar/decir cualquier cosa} ‘say anything possible’ or ‘say something random’. \textit{Cualquiera} under predicative verbs such as \textit{parecer} will be analysed as a degree adjective that has its origin in the FCI \textit{cualquiera} with the interpretation ‘any’.

{The basic data in this section is extracted from the \textit{Corpus del Español web dialects} (CdE\_web dialects, already introduced in \sectref{sec:1.5} and used in \chapref{sec:3}). The present study thus mainly focuses on diastratic variation and very little on social or diaphasic variation.}\footnote{{I asked about the use of} {\textit{cualquiera}} {in predicative position and under transitive verbs with the possibility of degree modification (e.g.} {\textit{es muy/re cualquiera}}{) in the online forum on Facebook} \href{https://www.facebook.com/groups/686284488187168/permalink/992126550936292/}{{{\textit{Lingüística Argentina}}}}, see \chapref{sec:1}. {18 native speakers responded to my questions. The informal questions asked are as follows: “How do you find the following sentences? Do they have the interpretation as represented in the brackets?” The following examples were provided:
  i) Me parece cualquier cosa. (para decir: me parece mal.) ‘I think \textit{cualquier cosa} (to say: I think it's wrong.)’
  ii) Hizo/Dijo cualquier cosa. (para decir: hizo/dijo algo malo) ‘He did/said \textit{cualquier cosa} (to say: he did/said something bad).’
  iii) La peli es cualquier cosa. (para decir: la pelí es mala). ‘The movie is cualquier cosa (to say: the movie is bad).’
  iv) Es tan/medio cualquier cosa. (para decir: Es tan/medio malo). It's so/half anything (to say: It's so/somewhat bad).
  v) Juan es más cualquiera que Marc. (para decir: Juan es peor que Marc.). ‘Juan is more \textit{cualquiera} than Marc (to say: Juan is worse than Marc).’
  vi) Esta peli es más cualquiera que la otra. (para decir: esta pelí es peor que la otra.). ‘This movie is more \textit{cualquiera} than the other one (to say: this movie is worse than the other one).’ Most speakers accepted} {\textit{cualquiera}} {in predicative position and under transitive verbs. However, not all degree adverbs can modify} {\textit{cualquiera}}{, as noted by native speakers. Most speakers accept} {\textit{re cualquiera}}{, in contrast to} {\textit{muy cualquiera}}. A systematic experimental study of speaker judgments will be performed in the future.} The quantitative analysis of the corpus results is represented as the ratio (in percentage) of various linguistic contexts and meanings of \textit{cualquiera}. The focus of my analysis is the distribution of EVAL and FCI of \textit{cualquiera} in Argentinian Spanish in the corpus data.

This section has the following structure. I will briefly recall the current state of research on \textit{cualquiera} in Argentinian Spanish from \chapref{sec:2} in \sectref{sec:4.1}. In \sectref{sec:4.2} I will show that the description in the literature is incomplete and very informal and will object to the suggestion that the source for the special use of \textit{cualquiera} in Argentinian Spanish should be derived from an analogy with the Italianism \textit{cualunque} ‘whichever’ (\cite{Rizzosalierno2013}). I will present new data on \textit{cualquiera} in \sectref{sec:4.3} that will serve my synchronic analysis in \sectref{sec:4.4}. And in \sectref{sec:4.5}, I will suggest that \textit{cualquiera} in Argentinian Spanish has lexicalised the speaker’s evaluation of the modal component denoted by \textit{cualquiera} as an FCI, which has triggered a reanalysis into an evaluative adjective or an evaluative noun. I will show further evidence for the analysis in \sectref{sec:4.6}.

\section{{Current} {state} {of} {research} {on \textit{cualquiera}} {in} {Argentinian} {Spanish}} \label{sec:4.1}

The item \textit{cualquiera} can have a pejorative or derogatory use (DER, a subtype of the evaluative use or function, EVAL) in Argentinian Spanish with copular predicates such as \textit{parecer} ‘to look like’ and \textit{ser} ‘to be’, similar to European Spanish, but, in addition, can be used predicatively and can have degree modification by degree adverbs (\citealt{Rizzosalierno2013}: example \REF{ex:key:2}, \sectref{sec:3.1}, \chapref{sec:2}). The following examples show degree modification of \textit{cualquiera} by the degree adverb \textit{completamente} ‘totally/completely’, and they are interpreted as ‘(very) bad’ or as ‘(complete) nonsense’. As I will argue in this chapter, the meaning ‘nonsense’ is a pejorative interpretation of ‘anything’ and thus another instance of DER:

\ea \label{ex:key:187} %ex 187
    \gll Es totalmente \textit{cualquiera} eso de haber puesto a Roger con la voz de Syd.\\
    is totally \textsc{cualquiera} that of having put to Roger with the voice  of  Syd\\
    \glt ‘Putting Roger with Syd's voice is total(ly) bullshit/nonsense/out of place/bad.’

    (\citealt{Rizzosalierno2013}: 28)
\z

\ea \label{ex:key:188} %ex 188
    \gll Eso de la tanga es completamente \textit{cualquiera}.\\
    that of the thong is completely \textsc{cualquiera}\\
    \glt ‘That thong thing is complete(ly) bullshit/nonsense/out of place/bad.’

    (\citealt{Rizzosalierno2013}: 28)
\z

According to \citet{Rizzosalierno2013}, the potential linguistic precedent of \textit{cualquiera} is possibly related to the Italianism \textit{cualunque}, which seems to have its origin in Lunfardo slang. \textit{Cualunque} in Argentinian Spanish (unlike Italian \textit{qualunque}, \citealt{Kellert2021a}) can also be used predicatively or attributively in a postnominal position and allows degree modification with \textit{muy} ‘very’ and other degree adverbs such as \textit{completamente}, as happens with ArgSp \textit{cualquiera} (\citealt{Kellert-francia-2023} for a historical overview and analysis of \textit{cualunque}).

\ea \label{ex:key:189-1} %ex 189
\gll Washizu se asustó porque sus zapatos staban sucios y su ropa era muy \textit{cualunque}.\\
    Washizu him scared because his shoes were dirty and his clothing was very \textsc{cualunque}\\
\glt ‘Washizu was scared because his shoes were dirty and his clothes were very {bad/out of place.}’\\
    (\citealt{Rizzosalierno2013}: 29)
\z


\ea \label{ex:key:190} %ex 190
    \gll Es una \textit{flaca} \textit{cualunque} que se cree que es Marilyn Monroe.\\
    is a woman \textsc{cualunque} that her believes that is Marilyn Monroe\\
    \glt ‘She’s an ordinary woman who thinks she’s Marilyn Monroe.’\\
    (\citealt{Rizzosalierno2013}: 29)
\z


Rizzo Salierno does not discuss cases of \textit{cualquiera} under transitive verbs, and he does not analyse the meaning ‘nonsense’, probably because he believes that \textit{cualquiera} has its origin in \textit{cualunque} and because he observed \textit{cualunque} only in predicative contexts. As Kellert \& Francia (in press) show \textit{cualunque} is indeed used mostly in predicative contexts or as a nominal modifier. They conclude that the transitive use of \textit{cualquiera} with the meaning ‘nonsense’ cannot have its origin in \textit{cualunque}. The transitive use of \textit{cualquiera} had been previously noted in the literature (see \sectref{sec:2.1.4}). \citet{Conde2011} assumes that the transitive use of \textit{cualquiera} as ‘nonsense’ is specific to Lunfardo. However, as of yet this use has not been investigated in depth.

\section{{Open} {questions} {and} {main} {hypothesis}} \label{sec:4.2}

The main question of this section is to determine the origin of the semantic and syntactic change of \textit{cualquiera} and to account for different meanings of \textit{cualquiera} as ‘ordinary’ and ‘nonsense’ and for their different syntactic functions as arguments of transitive verbs and as predicative elements.

Rizzo Salierno suggests that some analogy with \textit{cualunque} might be the cause of change. However, despite some similarities between \textit{cualunque} and \textit{cualquiera} in predicative contexts, they have a different distribution under transitive verbs (\citealt{Kellert-francia-2023}). Whereas \textit{cualquiera} occurs as a complement of verbs of saying, \textit{cualunque} does not:

\ea \label{ex:key:191} %ex 191
    \gll Me {contestó} / {dijo} \textit{cualquiera} / \textit{*cualunque}.\\
    me {answered} / {said} \textsc{cualquiera} / \textsc{cualunque}\\
    \glt ‘(S)he answered me with some bullshit/nonsense.’\\
    (\citealt{Rizzosalierno2013}: 26)
\z

Even if Rizzo Salierno was right about the assumption that \textit{cualquiera} has acquired its evaluative meaning from \textit{cualunque}, this assumption does not explain the origin of the evaluative meaning, which is absent from the Italian source item \textit{qualunque}. It just shifts the issue of how an indefinite can develop this type of evaluative meaning to \textit{cualunque}, but does not resolve it. Rizzo Salierno does not try to explain the origin of the evaluative use of ArgSp \textit{cualquiera} either historically or in the semantic structure of this item, and thus does not try to relate the evaluative function of \textit{cualquiera/cualunque} to the FCI meaning of these items.

I will fill this gap. Moreover, I will test Rizzo Salierno’s description of \textit{cualquiera} as a gradable adjective and investigate whether it behaves like a regular gradable adjective that can be modified by different degree modifiers (also including, e.g., items meaning ‘a bit’) and be used in comparative constructions (e.g. \textit{más cualquiera que} ‘worse than’). I will test this hypothesis with other degree modifiers and comparative constructions that are usually possible with degree adjectives.

As stated previously, my main hypothesis about the diachronic origin of the different readings of \textit{cualquiera} is that both the evaluative meaning ‘ordinary, bad’ (EVAL1) and the meaning ‘nonsense’ (EVAL2 or DER) have their origin in the free choice indefinite meaning ‘any’. I already showed in \chapref{sec:3} how this hypothesis applies to EVAL1, and I now show this for the special EVAL2 meaning found in Argentinian Spanish.

\ea \label{ex:key:192-1} %ex 192
    Free Choice \textit{cualquiera} ‘any’ > Adj \textit{cualquiera} ‘unremarkable/bad’ (EVAL1) or Noun \textit{cualquiera} ‘nonsense’ (EVAL2)
\z

In order to test this hypothesis for EVAL2, I will look at the distribution of \textit{cualquiera} in synchrony and diachrony in Argentinian Spanish in this section and offer an analysis that explains the change in \REF{ex:key:225}.

\section{{New} {data} {on} {\textit{cualquiera} }{in} {Argentinian} {Spanish}} \label{sec:4.3}

This section presents new data on \textit{cualquiera} in Argentinian Spanish, which will then be analysed in \sectref{sec:4.4}.

\subsection{{General} {distribution} {of} {\textit{cualquiera} }{in} {Argentinian} {Spanish}} \label{sec:4.3.1}

{The postverbal pronoun \textit{cualquiera} occurs 1.155 times in the corpus (per mil = 9.47). The special meaning/use of the pronoun \textit{cualquiera} in Argentinian Spanish ‘bullshit, nonsense’ occurs in object position as a bare pronoun under predicative verbs \textit{ser (177 occurrences)/parecer (37 occurrences)} and transitive verbs of \textit{hacer (2 occurrences), entender (6 occurrences, hablar (4 occurrences), mandar (23 occurrences), decir (9 occurrences)} out of 1.155 occurrences in total. The remainder of the occurrences has the FCI interpretation as in \textit{eliminar cualquiera de las opciones} ‘to eliminate any of the options’ (867 occurrences=75 percent out of 1.155 in total) (\tabref{tab:distribution-ArgSp.}). The subject function of the pronoun \textit{cualquiera} and the object function with a prepositional phrase inside the construction [\textit{cualquiera} de N] (e.g. \textit{eliminar cualquiera de las opciones} ‘to eliminate any of the options’) always has the interpretation of FCI.} 

\begin{table}
\begin{tabularx}{\textwidth}{QQQ}
\lsptoprule
 Linguistic context & EVAL & FCI ‘any' \\
 \midrule
 subject & -- & + \\
 \tablevspace
  \textit{cualquiera de N} & -- & + \\
 \tablevspace
 bare object pronoun & + (predicative verbs and verbs of saying in present and past tense) & + (non-indicative verb mood) \\
 \tablevspace
 frequency & 25 percent\newline (288 occ. out of 1.155) & 75 percent\newline (867 occ. out of 1.155) \\
 \lspbottomrule
\end{tabularx}
\caption{\label{tab:distribution-ArgSp.}Distribution of \textit{cualquiera} in Argentinian Spanish} %not in manuscript
\end{table}

{Here are some examples of the pronoun \textit{cualquiera} in the object position under indicative verbal mood of predicative verbs with the senses EVAL1 ‘bad’ and EVAL2 ‘bullshit, nonsense’.}

\ea \label{ex:key:4-1}
    \gll En resumen, la película es \textit{cualquiera}, si no fuera almodovar el que la dirige, diríamos que es una peli de cine clase b malísima.\\
    in summary, the movie is \textsc{cualquiera}, if \textsc{neg} were Almodóvar the that it directs, would.say that is a movie of cinema class b very.bad\\
    \glt ‘In short, the movie is bad, if it were not Almodóvar who directed it, we would say that it's a really bad B movie.' 
    
    (@ 276)
\z

\ea \label{ex:key:4-2}
    \gll es raro, es decir, ellos te piden que seas previsor pero después te piden hotel y fecha de viaje, es medio \textit{cualquiera} eso.\\
    is strange is to.say they you ask that be cautious but afterwards you ask hotel and date of trip is kind.of \textsc{cualquiera} that \\
    \glt ‘It's strange, that is, they ask you to be cautious but then they ask you to provide the hotel name and the date, that doesn't make much sense.' 
    
    (@277)  
\z

\ea \label{ex:key:4-3}
    \gll pero se nota que entendí \textit{cualquiera} jajaja\\
    but \textsc{refl} noticeable that understand \textsc{cualquiera} jajaja\\
    \glt ‘but it's obvious that I understood nonsense hahaha.’ 
    
    (@ 282)
\z

\ea \label{ex:key:4-4}
    \gll Le mandé \textit{cualquiera} jajajaja.\\
    to.them sent \textsc{cualquiera} jajajaja \\
    \glt ‘I told her nonsense hahaha.’  
    
    (@ 281)’ 
\z


\ea \label{ex:key:4-5}
     \gll Quiero que seas feliz. Estas haciendo \textit{cualquiera}.\\
    want.\textsc{1SG} that be.\textsc{2SG} happy be.\textsc{SG} doing \textsc{cualquiera}\\
    \glt ‘I want you to be happy. You are doing nonsense.’ 
    
    (@ 275)  
\z


\subsection{{Distribution} {of} {special} {uses} {of} {\textit{cualquiera} }{across} {Latin} {American} {varieties}} \label{sec:4.3.1-2} %section 4.3.2 in the document

In this subsection we will see that predicative \textit{cualquiera} and \textit{cualquiera} under transitive verbs such as \textit{mandar} ‘say’ are more frequent in Argentinian Spanish than in other varieties of Spanish.

The predicative \textit{cualquiera} as in \textit{es cualquiera} (\tabref{tab:key:8}) and \textit{parece cualquiera} (\tabref{tab:key:9}) is most frequent in Argentinian Spanish in contrast to other varieties (decimal numbers represent relative frequency in per million, i.e. how often a word or a word phrase occurs per one million words). The number 0.32 represents the mean of how many times \textit{es cualquiera} appears in all varieties per million words. The crucial observation is that the number of times \textit{es cualquiera} appears in Argentinian Spanish is three times higher than the mean. The number 1.05 represents the average number of instances of \textit{es cualquiera} per million in the Argentinian subcorpus, i.e. a set of one million words in the Argentinian subcorpus contains an average of 1.05 instances of \textit{es cualquiera}.


\begin{table}
% \begin{tabularx}{\textwidth}{ll rrrrrrrrrrrrrrrrrrrrrr}
% \lsptoprule
%  {{context}}          & {all} & {per} {million} & \multicolumn{2}{c}{Argentina} & \multicolumn{2}{c}{Bolivia} & \multicolumn{2}{c}{Chile} & \multicolumn{2}{c}{Colombia} & \multicolumn{2}{c}{Costa Rica} & \multicolumn{2}{c}{Cuba} & \multicolumn{2}{c}{Dom. Rep.} & \multicolumn{2}{c}{Ecuador} & \multicolumn{2}{c}{*Spain} & \multicolumn{2}{c}{Guatemala}\\
% \midrule
%  \textit{es cualquiera} &  633 & 0.32 & { 1.05} &   178 & { 0.20} &   8 & { 0.11} &   7 & { 0.20} &   33 & { 0.78} &   23 & { 0.38} &   24 & { 0.27} &   9 & { 0.21} &   11 & { 0.18} &   77 & { 0.24} &   13\\
%  \tablevspace
%   &  \multicolumn{2}{c}{Honduras}  &  \multicolumn{2}{c}{Mexico} & \multicolumn{2}{c}{Nicaragua}  &  \multicolumn{2}{c}{Panama}  &  \multicolumn{2}{c}{Peru} & \multicolumn{2}{c}{Puerto Rico} & \multicolumn{2}{c}{Paraguay} &  \multicolumn{2}{c}{El Salvador} &  \multicolumn{2}{c}{USA}  &  \multicolumn{2}{c}{Uruguay} & \multicolumn{2}{c}{Venezuela}\\
% \midrule
%   & { 0.43} &   15 & { 0.18} &   45 & { 0.25} &   8 & { 0.58} &   13 & { 0.24} &   26 & { 0.40} &   13 & { 0.27} &   8 & { 0.30} &   11 & { 0.36} &   59 & { 0.72} &   28 & { 0.24} &   24\\
% \lspbottomrule
% \end{tabularx}

\begin{tabularx}{\textwidth}{lrr Xrr}
    \lsptoprule
    Country & Absolute & Per million & Country & Absolute & Per million\\
\midrule
  \textbf{Argentina} &  \textbf{178} &  \textbf{1.05}     &    Honduras &  0.43 &   15      \\
  Bolivia &8 &  0.20       &    Mexico & 45& 0.18        \\
  Chile &   7 &  0.11           &    Nicaragua &    8 &  0.25    \\
  Colombia & 33       &  0.20 &      Panama &   13  & 0.58      \\
  Costa Rica &   23     & 0.78 &      Peru &   26 &  0.24        \\
  Cuba &    24 &  0.38          &    Puerto Rico &    13 & 0.40  \\
  Dom. Rep &   9       & 0.27 &      Paraguay &   8   & 0.27    \\
  Ecuador &  11        &  0.21 &     El Salvador &    11 & 0.30    \\
  Spain &  77          &   0.18 &   USA &     59  &  0.36       \\
  Guatemala &  13      &  0.24 &     Uruguay &  28   & 0.72      \\
            &        &          &    Venezuela &   24 & 0.24     \\
\midrule
           &           &       &  {total} &    {633} &  {0.32}           \\
  \lspbottomrule
\end{tabularx}
\caption{\label{tab:key:8} Distribution of \textit{es cualquiera} in CdE\_web dialects (highlighting my own).} %table 8
\end{table}

As we can see, the construction does not seem to be exluded in European Spanish, but seems to be extremely marginal there, as it also is in Chile, Bolivia, Colombia, and Mexico (with values ${\leq}$ 0.2 per million). It is more frequent in most other Latin American countries, but the value rises to more than 0.5 per million only in a small part of Central America (Panama, 0.58, and Costa Rica, 0.78) and in South America (in Uruguay, 0.72), whereas in Argentina it rises to 1.05 per million.

The same observation holds for \textit{parece cualquiera}, which is even more marginal than \textit{es cualquiera}, whose frequency is shown in \tabref{tab:key:10}. In most Latin American varieties in CdE, \textit{parece cualquiera} is not used at all. It is only used in Argentina, Spain, Mexico, and Peru. In Argentina, it is used with a strikingly higher frequency than in other countries (12 times higher than the average in the three other countries).


\begin{table}
% \begin{tabularx}{\textwidth}{XXXXXXXXXXXXX}
% \lsptoprule
%  {\textbf{context}} & \textbf{all} & \textbf{per} \textbf{million} & \textbf{Argentina} & \textbf{Bolivia} & \textbf{Chile} & \textbf{Colombia} & \textbf{Costa} \textbf{Rica} & \textbf{Cuba} & \textbf{Dom.} \textbf{Rep.} & \textbf{Ecuador} & \textbf{Spain} & \textbf{Guatemala}\\
%  { \textit{parece cualquiera}} & 30 & 0.02 & { 0.15}
%
%  25 &  &  &  &  &  &  &  & { 0.01}
%
%  3 & \\
% &  & \textbf{Honduras} & \textbf{Mexico} & \textbf{Nicaragua} & \textbf{Panama} & \textbf{Peru} & \textbf{Puerto} \textbf{Rico} & \textbf{Paraguay} & \textbf{El} \textbf{Salvador} & \textbf{USA} & \textbf{Uruguay} & \textbf{Venezuela}\\
% %\hhline%%replace by cmidrule{~~-----------} &  &  & { 0.01}
%
%  1 &  &  & { 0.01}
%
%  1 &  &  &  &  &  & \\
% %\hhline%%replace by cmidrule{~~-----------}
% \lspbottomrule
% \end{tabularx}

\begin{tabularx}{\textwidth}{lrr Xrr}
    \lsptoprule
    Country & Absolute & Per million & Country & Absolute & Per million\\
\midrule
  Argentina  &   25    &  0.15      &    Honduras    &  0     &  0    \\
  Bolivia    &   0    &  0      &    Mexico      &   1    &  0.01    \\
  Chile      &  0     &   0     &    Nicaragua   &  0     &   0   \\
  Colombia   &  0     &    0    &    Panama      &  0     &   0   \\
  Costa Rica &   0    &   0     &    Peru        &   1    &   0.01   \\
  Cuba       &  0     &    0    &    Puerto Rico &   0    &   0   \\
  Dom. Rep   &   0    &    0    &    Paraguay    &   0    &   0   \\
  Ecuador    &    0   &    0    &    El Salvador &    0   &   0   \\
  *Spain     &   3    &    0.01    &    USA         &   0    &  0    \\
  Guatemala  &    0   &    0    &    Uruguay     &  0     &  0    \\
             &  0     &    0    &    Venezuela   &   0    &   0   \\
\midrule
             &       &          &  {total} &  30    &  {0.02}     \\
  \lspbottomrule
\end{tabularx}

\caption{\label{tab:key:9} Distribution of \textit{parece cualquiera} in CdE\_web dialects (highlighting my own)}
\end{table}

The use of \textit{cualquiera} under the verb \textit{mandar} ‘say’ is also more frequent in Argentinian Spanish than in the other varieties as shown by the difference in percentages (75\% in Argentinian Spanish vs 25\% in all the other varieties).

Of the 44 occurrences of this construction in the corpus, 33 (= 75\%) are found for Argentina, as the \tabref{tab:key:10} shows. Moreover, the most frequent verb forms are the indicative present and past tense forms, which refer to some past or present speech events, i.e. someone said or is presently saying something \textit{wrong/nonsense} (\tabref{tab:key:10}).
%is this the correct table?

\begin{table}
\small
\begin{tabularx}{\textwidth}{l@{~~}r@{~~}r@{}YCCC@{}}
\lsptoprule
 {Context} & {All} & {Argentina} & {Other} {countries} & { {Indicative}}{ {Present}} { \mbox{Imperfect}}  {Perfect} & { {Subjunctive/}} \mbox{Imperative} & { {Gerund}} { {Infinitive}}\\
 \midrule
 \textit{mandan   cq.} & 8 & 6 & 2 & + & - & -\\
 \textit{mandaste cq.} & 5 & 3 & 2 & + & - & -\\
 \textit{mandaron cq.} & 5 & 5 & 0 & + & - & -\\
 \textit{mandé    cq.} & 4 & 4 & 0 & + & - & -\\
 \textit{mandar   cq.} & 3 & 2 & 1 & - & - & +\\
 \textit{mandando cq.} & 3 & 2 & 1 & - & - & +\\
 \textit{mando    cq.} & 3 & 3 & 0 & + & - & -\\
 \textit{manda    cq.} & 3 & 2 & 1 & + & - & -\\
 \textit{mande    cq.} & 3 & 1 & 2 & - & + & -\\
 \textit{mandó    cq.} & 2 & 2 & 0 & + & - & -\\
 \textit{mandaba  cq.} & 2 & 1 & 1 & + & - & -\\
 \textit{mandaban cq.} & 2 & 1 & 1 & + & - & -\\
 \textit{mandamos cq.} & 1 & 1 & 1 & + & - & -\\
\midrule
  {total} & { 100\%}   & { 75\%}   & { 25\%}  &  &  & \\
                &   (44) &   (33) &  (11) &  &  & \\
\lspbottomrule
\end{tabularx}
\caption{\label{tab:key:10} Distribution of lemma \textit{mandar cualquiera} in CdE\_web dialects} %Table 10
\end{table}

{Degree modification as in \textit{re/tan/muy/medio/totalmente/demasiado cualquiera} (‘very/a bit/fully/too non-sensical’) shows higher relative frequency in Argentina than in other countries}.   

\begin{table}[t]
\begin{tabularx}{.8\textwidth}{lYY}
    \lsptoprule
    Country & Abs. frequency & Per million \\
\midrule
%         ALL & 61 & 0.03 \\
    \textbf{Argentina} & \textbf{24} & \textbf{0.13} \\
    Bolivia & 2 & 0.05 \\
    Chile & 2 & 0.03 \\
    Colombia & 3 & 0.02 \\
    Costa Rica & 0 & 0.00 \\
    Cuba & 0 & 0.00 \\
    Dominican Republic & 0 & 0.00 \\
    Ecuador & 2 & 0.04 \\
    Spain & 8 & 0.02 \\
    Guatemala & 3 & 0.05 \\
    Honduras & 0 & 0.00 \\
    Mexico & 7 & 0.03 \\
    Nicaragua & 0 & 0.00 \\
    Panama & 0 & 0.00 \\
    Peru & 1 & 0.01 \\
    Puerto Rico & 1 & 0.03 \\
    Paraguay & 1 & 0.03 \\
    El Salvador & 1 & 0.03 \\
    United States & 4 & 0.02 \\
    Uruguay & 2 & 0.05 \\
    Venezuela & 0 & 0.00 \\
\midrule
    total & 61 & 0.03 \\
\lspbottomrule
\end{tabularx}
\caption{\label{tab:distribution-degree} Distribution of degree modification of \textit{cualquiera}}
\end{table}

\largerpage[1.5]
The difference in frequencies across Latin American varieties suggests that \textit{cualquiera} has undergone a change in Argentinian Spanish. Despite the clear difference in frequencies of the special uses between Argentinian and other Latin American varieties, the frequency of the special uses is very low in Argentinian Spanish compared to other phrases (e.g. \textit{decir algo} ‘say something’, with 408 occurrences, or \textit{es común} ‘it is common’, with 1,447 occurrences). One reasonable explanation of the low frequency of these uses is that they represent informal spoken speech (see also \citealt{Rizzosalierno2013}, \citealt{Conde2011}, among others).

The other Latin Spanish varieties will not be considered here, nor the question whether Argentinian Spanish has influenced the use of \textit{cualquiera} in other varieties.

\subsection{{Predicative} {\textit{cualquiera}}} \label{sec:4.3.2}

In this section, I will look into the predicative Argentinian Spanish \textit{cualquiera} in detail and see whether it behaves as a gradable adjective in any respect and whether it allows any kind of syntactic combination like usual gradable adjectives do (e.g. different combinations with adverbs, comparatives, etc.).

As already observed in \tabref{tab:key:4} and in \sectref{sec:4.1}, Argentinian Spanish predicative \textit{cualquiera} can be modified by degree adverbs that denote high degrees (i.e. adverbs meaning ‘very’ and the like). The most frequent modifying item is \textit{re}\footnote{{The detailed linguistic status of} {\textit{re}}{, as well as its origin, needs to be addressed in future research. In this book, I assume in line with \citet{Rizzosalierno2013} and \citet{KanwitTerán2020} that it is a degree adverb similar to} {\textit{muy}} {‘very’ that denotes a high degree. The form} {\textit{re}} {has been described as a derivative prefix from Latin that has grammaticalised as an intensifier (\citealt{KanwitTerán2020}: 257 and references therein). The intensifier} {\textit{re}} {has developed from a locative or spatial marker to a marker of repetition or iteration, and finally to an intensifier adverb (see \citealt{KanwitTerán2020}: 257 and references therein).}} (25 occurrences in the corpus). The most important observation is that the predicative \textit{cualquiera} has the meaning ‘nonsense’ if the subject of the predicative complement is a proposition or an event:

\ea \label{ex:key:193} %ex 193
    \gll ademas me parece \textit{re} \textit{cualquiera} eso de q se encajo una piña sola para hecharle la culpa a Mar\\
    moreover to.me seems very \textsc{cualquiera} that of that (s)he received a punch alone for throw.her the blame to Mar\\
    \glt ‘Moreover, that receiving a punch simply to blame Mar seems complete nonsense to me.’\\
    (@ 106)
\z

All occurrences of predicative \textit{cualquiera} in the Argentinain data of CdE\_web dialects co-occur with adverbial or verbal phrases that express whose negative evaluation it is that \textit{cualquiera} denotes (see the concept of “judge” in \citealt{Lasersohn2005}, among others; and \chapref{sec:3} on the judge parameter of \textit{tasty}{}-predicates). In almost all cases that I found in CdE\_web dialects, the evaluation is linked to the speaker, e.g. \textit{para mí} ‘for me’, \textit{me parece que} ‘It seems to me that…’. Of the 25 occurrences of \textit{parece cualquiera}, 17 are \textit{me parece}, two are \textit{te parece}, and one is \textit{parece}. All cases represent informal speech/dialogue:

\ea \label{ex:key:194}
    \gll No lo terminé de leer porque me pareció \textit{cualquiera}.\\
    \textsc{neg} it finish of read because me seemed \textsc{cualquiera}\\
    \glt ‘I did not finish reading it because it seemed \textit{nonsense} to me.’\\
    (@ 107)
\z


\ea \label{ex:key:195} %ex 195
    \gll Para mí, me parece \textit{cualquiera}.\\
    for me me seems \textsc{cualquiera}\\
    \glt ‘As for me, I find it \textit{nonsense}.’\\
    (@ 108)
\z 

\ea \label{ex:key:196} %ex 196
    \gll eso es \textit{cualquiera}, al menos para mí.\\
    that is \textsc{cualquiera} at least for me\\
    \glt ‘That is \textit{nonsense}, at least for me.’\\
    (@ 109)
\z

However, I also found a few cases of evaluative \textit{cualquiera} with a 3rd person opinion holder (see also \citealt{Lasersohn2005} for the same observation with \textit{tasty}{}-predicates):

\ea \label{ex:key:197} %ex 197
    \gll haría lo mismo pero si hay gente que no le gusta porque piensa que es \textit{cualquiera} (yo no pienso eso).....\\
    would.do the same but if there.are people that \textsc{neg} it like because think that is \textsc{cualquiera} (I \textsc{neg} think that)\\ 
    \glt ‘I would do the same but there are people that don’t like it because they think it is \textit{nonsense} (I don’t think that)….’

    (@ 110)
\z

In \sectref{sec:4.3.3}, we will see that \textit{cualquiera} with the meaning ‘nonsense’ also occurs under transitive verbs.

The meaning of \textit{cualquiera} as ‘nonsense’ can be reinforced by degree adverbs that denote high degrees such as \textit{muy} ‘very’, \textit{demasiado} ‘too much, very’, \textit{tan} ‘such’, \textit{re} ‘very’. Degree modification of \textit{cualquiera} as ‘nonsense’ is unexpected under the assumption that \textit{cualquiera} is a noun when it means ‘nonsense’. I will return to this puzzle in my analysis in \sectref{sec:4.4}. I translate in this case the modified \textit{cualquiera} as ‘quite nonsensical’, ‘such utter nonsense’, or as ‘completely out of place’:

\ea \label{ex:key:198} %ex 198
 \gll Me parece todo tan \textit{cualquiera}.\\
to.me seems everything so \textsc{cualquiera}\\
\glt ‘Everything seems such utter nonsense to me.’

(@ 111)
\z

If the subject of the predicative \textit{cualquiera} refers to a person or an entity, the meaning of \textit{cualquiera} can take on different flavours already discussed in \chapref{sec:3}, such as ‘ordinary, unremarkable’ and the pejorative meaning of ‘bad’.

\ea \label{ex:key:199} %ex 199
  \gll Yo, que soy tan \textit{cualquiera} como \textit{cualquiera}.\\
I that am so \textsc{cualquiera} as \textsc{cualquiera}\\
\glt ‘I who am as ordinary as anyone else.’\footnote{This \textit{cualquiera} is a pronoun (see \chapref{sec:2}).}

    (@ 112)
\z

Argentinian Spanish also has \textsc{un} N \textit{cualquiera} in predicative position just as in European Spanish (e.g. \textit{una fiesta re cualquiera} ‘a very bad party’). The big difference to European Spanish is that \textit{cualquiera} can also be modified by a degree adjective and be used as a predicative adjective as in \REF{bkm:Ref63461360}.

Note that \textit{cualquiera} can predicate over plural entities as well, as shown in \REF{bkm:Ref150786385}. In the following context, \textit{cualquiera} expresses an evaluative property of entitities such as pictures according to a judge (here: B). In the following context, the evaluation concerns what the pictures represent (i.e. namely that they show a bad example):

\ea 
\label{bkm:Ref150786385} %ex 200 
\ea \gll A: Esas fotos que publican parece que dan mal ejemplo.\\
    { } those pics that publish seem that give bad example\\
\glt A: ‘It seems that these pics that they publish give the wrong idea.’\\
\ex \gll B: Sí, son re \textit{cualquiera}.\\
    {  } yes are very \textsc{cualquiera}\\
\glt B: ‘Yes, they [these photos] are completely out of place. The idea that these pictures represent does not make sense.’\\
  (@ 113)
\z
\z 

The following data show that \textit{cualquiera} need not necessarily combine with adverbs that denote high degrees (against \citeauthor{Salierno2013}’s 2013 observation). Actually, we find two occurences with modification by the adverb \textit{medio}, which literally denotes ‘half’ and means ‘a bit/almost’ when combined with gradable adjectives as in \REF{bkm:Ref150786427} and \REF{bkm:Ref150786436} and \textit{cualquiera} as in \REF{bkm:Ref150786444} and \REF{bkm:Ref150786457}:

\ea \label{bkm:Ref150786427}  %ex 201
\gll \textit{Es} \textit{medio} \textit{complicado} este asunto.\\
    is half complicated this case\\
\glt ‘It’s a bit complicated, this case.’\\
    (@ 114)
\z


\ea \label{bkm:Ref150786436}  %ex 202
    \gll matemáticamente \textit{es} \textit{medio} \textit{imposible} que todas las personas puedan pasar por esa situación\\
    mathematically is half impossible that all the people can go through that situation\\
    \glt ‘Mathematically it’s almost impossible for all the people to go through that situation.’\\
    (@ 115)
\z


\ea \label{bkm:Ref150786444} %ex 203
\gll es que la defensa del rugby por los valores que entraña \textit{es} \textit{medio} \textit{cualquiera.} O sea, a mí es un deporte que no me atrae.\\
    is that the defence of rugby of the values that involves is half \textsc{cualquiera} that is to me is a sport that not me attracts\\
    \glt ‘Defending rugby for the values it involves is a bit nonsensical/out of place. That is, it’s a sport that does not appeal me.’\\
    (@ 116)
\z


\ea \label{bkm:Ref150786457}  %ex 204
\gll \textit{es} \textit{medio} \textit{cualquiera} eso\\
    is half \textsc{cualquiera} that\\
\glt ‘That is a bit \textit{nonsensical}.’\\
    (@ 117)
\z

I even found that the diminuitive suffix \textit{-ita} can combine with \textit{cualquiera}, and in this case the speaker attenuates the bad quality ‘a bit nonsensical’:

\ea \label{ex:key:205} %ex 205
    \gll Bueno, cerramos este post que \textit{es}  \textit{cualquierita}.\\
    well close this post that is \textsc{cualquiera}.\textsc{dim}\\
    \glt ‘Well, let’s close this post because it’s a bit nonsensical.’\\
    (@ 118)
\z

If \textit{cualquiera} has been reanalysed as a gradable adjective, as such data suggest, it should be possible to find it in comparative clauses such as ‘he is worse than anyone else’. Surprisingly, the comparative with \textit{más que} was not found in CdE\_web dialects. However, the following case from the corpus might be interpreted as an elliptical comparative clause:

\ea \label{ex:key:206} %ex 206
    \gll el reggae es \textit{cualquiera}. Mas [cualquiera] el reggae nacional.\\
    the reggae is \textsc{cualquiera} more \textsc{[cualquiera]} the reggae national\\
    \glt ‘Reggae is \textit{bad}. National reggae is (even) worse.’\\
    (@ 119, interpolation my own)
\z

There are some examples of comparatives with \textit{cualquiera} on the Internet:

\ea \label{ex:key:207} %ex 207
    \gll Más \textit{cualquiera} que el otro.\\
    more \textsc{cualquiera} than the other\\
    \glt ‘worse than the other one.’\\
    (@ 120)
\z

\ea \label{ex:key:208} %ex 208 
    \gll Y con los libros me pasó más o menos lo mismo aunque no los pienso seguir porque son más \textit{cualquiera} que la serie por ahí.\\
    and with the books me happened more or less the same although \textsc{neg} them think follow because are more \textsc{cualquiera} than the series for there\\
    \glt ‘And with the books more or less the same thing happened to me, although I will definitely not follow them because they are worse than the series that is out there.’\\
    (@ 121)
\z

\ea \label{ex:key:209} %ex 209
    \gll me parece mas \textit{cualquiera} que la cumbia\\
    to.me seems more \textsc{cualquiera} than the cumbia\\
    \glt ‘It seems worse to me than the cumbia.’\\
    (@ 122)
\z


It thus seems that \textit{cualquiera} does occur in comparative structures, although very rarely, so the corpus is probably too small to detect it or \textit{cualquiera} is still on its way of being reanalysed as a gradable adjective. Another argument that \textit{cualquiera} has been already reanalysed as a gradable adjective is shown by the difference in interpretation of \textit{cualquiera} under negation. Recall from \chapref{sec:2} that, when it is used as an FCI, \textit{cualquiera} has a ‘not just any’ interpretation under negation. However, \textit{cualquiera} has a different interpretation when it occurs in predicative position as an adjective in \REF{bkm:Ref63460902}, in contrast to the indefinite pronoun in \REF{bkm:Ref63460922}.

\ea \label{bkm:Ref63460902}  %ex 210
    \gll No es tan \textit{cualquiera} pero sí medio \textit{cualquiera}.\\
    \textsc{neg} is so \textsc{cualquiera} but yes half  \textsc{cualquiera}\\
    \glt ‘It’s not that \textit{bad} but it is a bit bad.’\\
    (@ 123)
\z

\ea \label{bkm:Ref63460922}  %ex 211
\gll es un portal, y no es \textit{cualquiera} es el mejor portal de autos del país\\
    is a platform and \textsc{neg} is \textsc{cualquiera} is the best platform of cars in.the country\\
\glt ‘It’s a platform and not just any platform. It’s the best car platform in this country.’\\
    (@ 124)
\z

This section has shown that \textit{cualquiera} is used as a gradable category with evaluative meanings auch as ‘nothing special’, ‘bad’, ‘unremarkable’, ‘nonsensical’, or ‘out of place’ in predicative contexts in Argentinian Spanish. It can be modified with adverbs \textit{re/muy/tan/medio} ‘very/a bit’ and to a lesser extent with comparative adverbs such as \textit{más} ‘more’.

\subsection{{Pronominal} {\textit{cualquiera}}{ }{under} {transitive} {verbs}} \label{sec:4.3.3}

Until now in this chapter, I have presented data from Argentinian Spanish on predicative \textit{cualquiera} in which the item at issue has a series of evaluative meanings such as ‘bad’ or ‘nonsensical’.

In sentences with imperative mood, which I also analyse as a modalised constructions (see \citealt{am2018} and \sectref{sec:3.2.3}), ArgSp \textit{cualquiera} is interpreted as an indefinite pronoun ‘anything’ (which would correspond to \textit{cualquier cosa} in European Spanish). Recall from \sectref{sec:2.1.4} that \textit{cualquiera} can be used as a substitute for \textit{cualquier cosa} in European Spanish:

\ea \label{ex:key:212} %ex 212
    \gll {\textit{¡}}\textit{Hagamos} \textit{cualquiera!} (hacer algo, hacer lo que sea, o como que da lo mismo hacer una cosa u otra).\\
do \textsc{cualquiera} do something do it what is or as that give the same do one thing or another\\ 
    \glt ‘Let’s do something! (do something, do anything, or it doesn’t matter whether it’s this thing or the other).’

(@ 125)
\z

In the following two examples \textit{digo cualquiera} and \textit{hace cualquiera} can be interpreted as ‘I say anything possible’ and ‘she does anything possible’, respectively, i.e. \textit{cualquiera} can be interpreted as a Free Choice Indefinite pronoun ‘anything’ under covert subtrigging. But another possible interpretation is that the speaker uses \textit{cualquiera} to qualify or evaluate the result state of the verbs ‘say’ and ‘do’ as ‘nonsense’ (EVAL2):

\ea \label{ex:key:213} %ex 213
\gll al rato volvió porque sabe que yo dormida \textit{digo} \textit{cualquiera}\\
   after while came.back because know that I asleep say \textsc{cualquiera}\\
\ea ‘She came back because she knows that when I sleep, I say anything possible.’ (FCI)\\

\ex ‘She came back because she knows that when I sleep, I talk nonsense/rubbish.’ (EVAL2)\\
    (@ 126)
\z 
\z 

\ea \label{ex:key:214} %ex 214
\gll se va cuando quiere, hace \textit{cualquiera}.\\
    her/him goes.out when wants does \textsc{cualquiera}\\
\ea  ‘(S)he goes out when (s)he wants to do so, (s)he does anything she wants. (FCI)\\
\ex ‘(S)he goes out when (s)he wants to do so, (s)he does nonsense.’ (EVAL2)\\
    (@ 127)
\z
\z 

\largerpage
It thus seems that that the interpretation of \textit{cualquiera} under certain transitive verbs in the indicative as ‘anything’ (FCI) or ‘nonsense’ (EVAL2) is context-dependent in Argentinian Spanish. In contrast, if \textit{cualquiera} occurs under the verb \textit{mandar} ‘to say’, it has only the evaluative meaning ‘to talk nonsense’, as has already been observed in \sectref{sec:2.1.4}. An example is given in \REF{bkm:Ref150786556}.

\ea \label{bkm:Ref150786556} %ex 215
\gll No supe qué decirle, así que le mandé \textit{cualquiera}.\\
not knew what say.her so that her told \textsc{cualquiera}\\
\glt ‘I did not know what to tell her, so I told her nonsense.’\\
    (@ 128)
\z

A crucial observation is that ArgSp \textit{cualquiera} also has an evaluative meaning like ‘nonsense’ when \textit{mandar} appears under a modal verb:\footnote{The context of the example is: \textit{Aca el problema es que estos tipos, como los chantas KK, se abusan porque hay fanáticos como este Juan emilio que no ven la realidad, viven en otro mundo y les pueden mandar cualquiera total no les da la cabeza para darse cuenta que lo estan verseando.}
‘The problem here is that these guys, like the KK charlatans, hurt themselves because they are fanatics like this Julio Emilio who do not see reality, they live in a different world and you can tell them complete bullshit, they don’t have the intelligence to realise that people are mocking them.’}

\ea \label{ex:key:216} %ex 216
    \gll y les pueden mandar \textit{cualquiera} total\\
    and that her sent \textsc{cualquiera} \textsc{total}\\
    \glt ‘and you can tell them complete bullshit.’\\
    (@ 129a)
\z

The evaluative reading here is reinforced by the degree adverb \textit{total} ‘complete’. This observation suggests that \textit{cualquiera} under \textit{mandar} ‘to say’, no matter what mood \textit{mandar} has, always has the meaning of ‘to talk nonsense/bullshit’ in contrast to other verbs that have different interpretations depending on the verbal mood. The only thinkable explanation is that the evaluative meaning of \textit{cualquiera} as ‘nonsense’ has been lexicalised in the construction of \textit{mandar cualquiera}, whereas with other verbs, the interpretation of \textit{cualquiera} depends on the verbal mood.

\subsection{{Summary} {of} {the} {data}} \label{sec:4.3.4}

Under modal verbs or in modalised constructions such as subjunctives and imperatives, ArgSp \textit{cualquiera} has an FCI function in the sense of ‘anything’. Non-modal contexts such as transitive verbs in the indicative past and present tense may also trigger FCI if \textit{cualquiera} is licensed under subtrigging. \textit{Cualquiera} under transitive verbs is something typical for Argentinian Spanish but not for other varieties of Spanish. However, in certain contexts that express evaluation and with certain verbs such as \textit{mandar, cualquiera} has an evaluative interpretation of ‘nonsense/bullshit’. In predicative contexts such as with copular verbs, \textit{cualquiera} has only the evaluative interpretation in Argentinian Spanish. It is in this context that it can be modified by different degree adverbs and occur in degree comparatives, although modification with degree comparative adverbs occurs less frequently than modification with simple degree adverbs such as ‘very/a bit’. The interpretation of \textit{cualquiera} as ‘ordinary’, ‘nonsense’, or ‘anything’ depends strongly on the grammatical category. If it is a pronoun, it has the interpretation of ‘anything’. If it is a noun, it has the interpretation of ‘nonsense’, and if it is an adjective, it has the interpretation of ‘ordinary’ or ‘nonsensical’.


\begin{table}
\fittable{
\begin{tabular}{lllll}
\lsptoprule
%\hhline%%replace by cmidrule{~----}
 {} & {Indicative} {present} {or} {past} {tense} & {Modal} & {\textit{mandar} }{‘say’} & {Predicative}\\
  \midrule
 {FCI} & + (subtrigging) & + & – & –\\
 {EVAL} & + & – & + & +\\
\lspbottomrule
\end{tabular}
}
\caption{\label{tab:key:11} Different interpretations in relation with verbal mood and verb type (+/– predicative verb)}
\end{table} %Table 11

\section{{Analysis} {of} {different} {uses} {of} {\textit{cualquiera}}{ }{in} {Modern} {Argentinian} {Spanish}} \label{sec:4.4}

In this section I will give an analysis of the different syntactic functions of \textit{cualquiera} in Modern Argentinian Spanish and discuss their origin. I will show that the indefinite use of \textit{cualquiera} ‘any’ is the origin of the evaluative functions of \textit{cualquiera} ‘ordinary’ and ‘nonsense’.

The data with degree modification suggests that ArgSp \textit{cualquiera} used in predicative contexts is the result of a reanalysis as a gradable adjective denoting a relation between degrees and individuals. It thus behaves in a way similar to English \textit{tall}: \textit{Juan is very tall} denotes that ‘Juan is tall to degree d, and degree d denotes a high degree on a scale of tallness’. \textit{Cualquiera} is thus a predicate denoting a scale of ordinary qualities or the pejorative version of it (i.e. bad qualities; see \sectref{sec:3} on the pejoration of ‘common’, ‘ordinary’). This kind of predicate has a subjective value, i.e. the badness is not objectively measurable in contrast to tallness. The gradable predicate thus takes an additional argument, which is the experiencer of the bad quality (see \sectref{sec:4.3}):

\ea \label{ex:key:217}
\ea \label{ex:key:217a} 
\gll la peli es re \textit{cualquiera}.\\
the movie is very \textsc{cualquiera}\\
\glt ‘The movie is very ordinary.’

\ex \label{ex:key:217b}
    Adj° \textit{cualquiera/cualunque}\footnote{{However, both} {\textit{cualquiera}} {and} {\textit{cualunque}} {are invariable adjectives as they do not inflect for person and gender. Number variation is possible in ArgSp} {\textit{las chicas son cualquieras/cualunques}} {‘the girls are nothing special’ (see also \citealt{Rainer1993} and \sectref{sec:2.1}).}} = \textit{λd.λx}. bad/ordinary(\textit{d})(\textit{x}) for some experiencer E

\z 
\z

This kind of predicate has a subjective value. Thus, cualquiera in \REF{ex:key:217a} would be interpreted as \REF{ex:key:217b}.

\ea \label{ex:key:218}
    [\textsubscript{DegP} re/medio/más/demasiado/tan [\textsubscript{AdjP} cualquiera/cualunque]]
\z

As I showed in \sectref{sec:4.3.2}, \textit{cualquiera} can also have the meaning of ‘nonsense’ when the object of evaluation is not an entity or person, but a proposition or an event, as in \REF{ex:key:219}:

\ea \label{ex:key:219}
    \gll Lo que me decís {es} \textit{cualquiera}.\\
    the.N that me say is \textsc{cualquiera}\\
    \glt ‘What you are telling me is nonsense.’\\
    (@ 130)
\z

The difference in meaning between ‘nonsense’ and ‘ordinary’ or the pejorative version of the latter meaning as ‘bad’ is related to the syntactic and semantic type of the syntactic subject of the predicative \textit{cualquiera.} In case the subject is a sentence or a proposition \textit{p} of semantic type <s,t> (i.e. the characteristic function of a set of possible worlds), \textit{cualquiera} means ‘nonsense’. In the following example, \textit{cualquiera} expresses evaluation with respect to what has been said, i.e. the argument of the verb of \textit{decir} ‘to say’. This argument is a proposition. In this example, \textit{cualquiera} is a bare noun, precisely the same one that occurs as a noun under verbs such as verbs of saying, such as ArgSp \textit{decir/mandar} ‘talk nonsense’, that usually take propositions as their arguments (e.g. \textit{decir/mandar}) (see \sectref{sec:4.5.2} for details).

\ea \label{ex:key:220}
    λp. cualquiera (p) = ‘nonsense’

\textit{decir/mandar cualquiera} = ‘talk nonsense’
\z

One important issue that needs to be addressed by this analysis is the fact that in Spanish and Romance languages in general, bare singular nouns are rare in contrast to indefinite singular nouns with an indefinite article (e.g. \textit{decir una tontería} ‘say a stupid thing’). This restriction does not apply to plurals (e.g. \textit{decir tonterías} ‘say stupid things’). I assume that the reason why \textit{cualquiera} appears as a bare noun rather than as an indefinite noun (*\textit{decir un/una cualquiera}) or a plural noun (*\textit{decir cualquieras}) is due to its pronominal origin. The pronoun \textit{cualquiera} ‘anything’ is invariable for plural or singular. Moreover, \textit{cualquiera} that takes a proposition as its argument is described as a kind property; that is, it describes the kind of things that one can say (see \sectref{sec:4.5.2} for details). Recall that \textit{cualquiera} that predicates over propositions can also appear with adverbs that reinforce the evaluative property as ‘utter nonsense’ or ‘very nonsensical’. In this case the usual modification is a nominal modifier \textit{total} ‘complete’ as repeated here:

\ea \label{ex:key:221} %ex 221
    \gll y  les  pueden  mandar  \textit{cualquiera} total\\
    and  that  her  sent  \textsc{cualquiera} {total}\\
    \glt ‘and you can tell them complete bullshit.’\\
    (@ 129b)
\z

However, there are also cases of \textit{cualquiera} with the meaning ‘nonsense’ and degree modification with (\textit{tan, re} ‘very’) when \textit{cualquiera} is used in a predicative position, as already shown in \REF{bkm:Ref150786811} and repeated here.

\ea \label{ex:key:222}
\gll Me parece todo tan \textit{cualquiera} escribirlo casi me da vergüenza.\\
to.me seems everything so \textsc{cualquiera} write.it almost to.me gives shame\\ 
\glt ‘Everything seems such utter nonsense to me writing it down. It’s almost embarrassing to me.’

(@ 111)
\z

This shows that \textit{cualquiera} can change its syntactic status as an adjective or a bare noun depending on the syntactic context. This flexibility in word categorisation is expected under my semantic analysis, as \textit{cualquiera} denotes a property and whether this property is represented as a noun or as an adjective depends on the syntax.

The suggested analysis can account for the difference between \textit{cualquiera} and \textit{cualunque} in Argentinian Spanish. The element \textit{cualunque} does not have the meaning of ‘nonsense’, as it cannot be used to predicate over propositions and thus be used under transitive verbs such as verbs of saying (see \sectref{sec:4.2}). The difference in meanings between ‘bad/ordinary’ and ‘nonsense’ is a further argument against deriving \textit{cualquiera} from the analogy of \textit{cualunque} (against \citealt{Rizzosalierno2013}). I believe that both items \textit{cualquiera} and \textit{cualunque} have developed the gradable property in a parallel way as a result of reanalysis of FCIs as (gradable) predicates.

\section{From FC use to evaluative meanings in Argentinian Spanish}
\label{sec:4.5}

In this section, I will provide an analysis of it \textit{cualquiera} in Argentinian Spanish and show how it changed to the evaluative \textit{cualquiera} analysed in \sectref{sec:4.5.1}. I will suggest different diachronic paths of \textit{cualquiera} as ‘ordinary’ and as ‘nonsense’. This suggestion will be tested on diachronic data of Argentinian Spanish in \sectref{sec:4.6.3}.

\subsection{{The} {change} {from} {the} {postnominal} {modifier} {to} {the} {adjective} {\textit{cualquiera}} {in} {Argentinian Spanish}} \label{sec:4.5.1}

I assume that the element \textit{cualquiera} with the meaning ‘ordinary’ or ‘bad’ (EVAL1) is diachronically derived from \textit{cualquiera} in \textsc{un} N \textit{cualquiera} by lexicalising the evaluative property of ‘ordinary, bad’ and by being used as a predicative adjective without a noun:

\ea{ \label{ex:key:223} %ex 223
\gll Juan es un hombre \textit{cualquiera}. \\
Juan is a man \textsc{cualquiera}\\}\jambox*{(EurSp and ArgSp)}
\glt ‘Juan is a man like any other man’\jambox*{SIM}

\glt ‘Juan is an ordinary man’\jambox*{COM}

STAGE 1 = evaluative use of the indefinite \textit{cualquiera}
\z 

\ea{ \label{ex:key:224} %ex 224
\gll Juan es un hombre muy \textit{cualquiera}. \\
Juan is a man very \textsc{cualquiera}\\}\jambox*{(ArgSp)}
\glt ‘Juan is a (very) ordinary man.’ 

Lexicalisation of COM and reanalysis into an adjective

STAGE 2 A= reanalysis into an adjective that modifies a noun.
\z 

\ea{ \label{ex:key:225} %ex 225
\gll Juan es (muy) \textit{cualquiera} \\
Juan is (very) \textsc{cualquiera}\\}\jambox*{(ArgSp)}
\glt ‘Juan is (very) ordinary.’ 

Separation from the noun

STAGE 2 B= reanalysis into an adjective without nominal modification.
\z

The diachronic steps can be formalised as follows. \textit{Cualquiera} passed from being analysed as a property of individuals or kinds of individuals into a property of degrees that can be used with or without a noun:

\ea \label{ex:key:226} %ex 226

\ea Stage 1

     λ\textit{P}.λ\textit{x\textsubscript{} }[(cualquiera (\textit{P}))(\textit{x})]

‘property (of a kind of) individuals’    

\ex Stage 2 A

    λ\textit{d}. λ\textit{P}. λ\textit{x}. [(cualquiera (P)(\textit{d}))(\textit{x})]

‘gradable property of individuals realised as a nominal adjective’  

\ex Stage 2 B

    λ\textit{d}. λ\textit{x} [(cualquiera (\textit{d}))(\textit{x})]

‘gradable property of individuals realised as a predicative adjective’     
\z
\z

The empirical evidence of these diachronic steps will be discussed in \sectref{sec:4.6.3} and in \chapref{sec:5}.
I assume that the noun \textit{cualquiera} took a different diachronic path.

\subsection{{The} {change} {from} {a} {pronoun} {into} {a} {noun} {\textit{cualquiera}} {‘nonsense’} {in} {Argentinian Spanish}} \label{sec:4.5.2}

My hypothesis is that \textit{cualquiera} under transitive verbs with the evaluative meaning ‘nonsense’ (EVAL2) (see \sectref{sec:4.4}) has its origin in the indefinite pronoun \textit{cualquiera} with the quantificational force of ‘anything’. Recall from \sectref{sec:2.1.4} that \textit{cualquiera} can be used as \textit{cualquier cosa} in ArgSp:

\ea \label{ex:key:227} %ex 227
    cualquiera\textsubscript{transitive} ‘any’ under ‘say’ = λP. λQ. Ǝx[Q(\textit{x}) \& P(x)] 
    
    ‘anything’

FC interpretation = everything is a possible option.
\z

Recall from \sectref{sec:2.1} that in European Spanish, inanimate \textit{cualquiera} requires mentioning the domain in the prior context. As the following example from European Spanish shows, \textit{cualquiera} stands for \textit{cualquier interpretación} (see \sectref{sec:2.1}).

\ea \label{ex:key:228} %ex 228
   \gll Aunque puedan admitir un alto número de interpretaciones (pero no admiten \textit{cualquiera}). Hay también una tercera posibilidad.\\
    although can admit a high number of interpretations but not admit \textsc{cualquiera} there.is too a third possibility\\
    \glt ‘Even though they can admit a high number of interpretations (not not just any). There is also a third possibility.’\\
    (@ 131; EurSp)
\z

In Argentinian Spanish, \textit{cualquiera} can be used to refer back to an item of the a preceding context. In the following examples, the preceding context has \textit{cualquier cosa}, and \textit{casi cualquiera} refers to \textit{cualquier cosa}. Both forms are the argument of the verb \textit{necesitar} ‘to need’ in  \REF{ex:key:229} and of \textit{ocurrir} ‘to happen’. Note also that \textit{casi} usually modifies FCIs (see \chapref{sec:3} for modification with \textit{casi} ‘almost’):

\ea \label{ex:key:229} %ex 229 \label{bkm:Ref150786782} 
\gll \textit{Cualquier} \textit{cosa} que necesiten (bueno, \textit{casi} \textit{cualquiera!}) nos avisan!\\
    \textsc{cualquier} thing that need well almost \textsc{cualquiera} us inform\\
    \glt ‘If you need anything (well, almost anything!) tell us!!’\\
    (@ 132)
\z

\ea \label{bkm:Ref150786875} %ex 230 
    \gll Hoy en día, uno se imagina que podía estar ocurriendo \textit{casi} \textit{cualquier} \textit{cosa.} \textit{Casi} \textit{cualquiera} menos, probablemente, lo que ocurría.\\
    today in day one her/him imagines that could be happening almost \textsc{cualquier} thing almost \textsc{cualquiera} minus probably the that happened\\
    \glt ‘Nowadays, one imagines that anything could be happening. Anything except what was actually happening, probably.’\\
    (@ 133)
\z

As already noted in \sectref{sec:2.1.4}, \textit{cualquiera} in Argentinian Spanish, but not in European Spanish, can be used as \textit{cualquier cosa} without previous mention of \textit{cualquier cosa}. \textit{Cualquiera} has lexicalised the meaning of \textit{cualquier cosa} in Argentinian Spanish and does not require previous mention of \textit{cualquier cosa}:

\ea \label{bkm:Ref150786811} %ex 231

Respondí/Entendí/Hice/Dije [\textsubscript{DP} \textit{cualquiera}]

    = ‘I answered/understood/did/said anything.’
\z

There is a pragmatic reason why \textit{cualquiera} as a verbal argument has the meaning of \textit{cualquier cosa} and not a more specific meaning, such as \textit{cualquier libro} ‘any book’. \textit{Cualquiera} in object position without any previous discourse can only mean \textit{cualquier cosa}, because \textit{cualquier cosa} is more generic than \textit{cualquier libro}. A similar case can be demonstrated with elliptical nouns that appear without previous mention. For example, \textit{me gusta comer} ‘I like eating’ is interpreted as ‘I like eating food or edible things’ and not as ‘I like eating a specific thing, such as chocolate.’


Two questions arise with respect to \REF{bkm:Ref150786811}. The first is what licenses the use of FCI \textit{cualquiera}, or \textit{cualquier cosa} for that matter, under perfective verbs in \REF{bkm:Ref150786875}, given that the licensing of FCIs usually happens under modal operators rather than under perfective verbs (see \chapref{sec:2}). The second question is how to account for the evaluative reading of \textit{cualquiera} under transitive verbs, e.g. \textit{decir}/\textit{mandar cualquiera} ‘to talk nonsense’ or \textit{hacer cualquiera} ‘to do something out of place/nonsense’ (see \sectref{sec:4.3.3}).

To address the first question, I assume that FCI \textit{cualquiera} in episodic contexts such as under transitive perfective verbs (e.g. \REF{ex:key:105} in \sectref{sec:2.3} repeated in \REF{bkm:Ref63461360}) can be licensed under subtrigging where the modal is contained inside a covert relative clause, e.g. \textit{que podría responder} ‘that I could possibly answer’ (see \sectref{sec:2.2} for subtrigging and the formal representation of it):

\ea \label{bkm:Ref63461360}  %ex 232
\gll Respondí \textit{cualquiera}  [que podría decir/responder].\\
responded \textsc{cualquiera} [that could say/respond]\\
    \glt ‘I answered anything (I could answer/say).’
\z

Another way in which the Argentinian Spanish pronominal FCI \textit{cualquiera} ‘anything’ can be licensed in episodic contexts is to assume a covert modal inside the matrix clause, as I represent informally in \REF{bkm:Ref63461220}, e.g. an agent’s decision-based modal that triggers the random choice interpretation (see \sectref{sec:2.2}):

\ea \label{bkm:Ref63461220} %ex 233 
\gll Respondí \textit{cualquiera}.\\
responded \textsc{cualquiera}\\
    \glt ‘I answered something random.’
\ea ‘I said something and this something fulfills my decision to say any possible thing.’
(Agent’s decision-based modality)
\ex ‘The thing I actually said could have been anything else than what I actually said.’ (Counterfactual modal)

    (@ 33)
\z
\z

The two licensing conditions (modal operator inside a relative clause, i.e. subtrigging, or inside the matrix clause as random choice modality) create two different truth conditions. Under subtrigging in \REF{bkm:Ref63461360}, \textit{cualquiera} refers to many speech events in the actual world (I said a, b, c, d….) and under random choice in \REF{bkm:Ref63461220}, \textit{cualquiera} refers to a single speech event in the actual world and the actual speech event is contrasted to other speech events that the speaker could have said or that agree with agent’s decision. There is evidence that we need both analyses, i.e. licensing of \textit{cualquiera} under subtrigging and under some modal that triggers random choice. In cases where \textit{cualquiera} can be modified with adverbs like \textit{casi} ‘almost’ or with \textit{except}{}-phrases (see \sectref{sec:2.1}), \textit{cualquiera} is licensed under subtrigging. But when \textit{cualquiera} has an existential inference in a particular context, it is licensed under a modal that triggers random choice. I will first derive the evaluative meaning of \textit{cualquiera} from FCIs under subtrigging, as shown in \REF{bkm:Ref63461360}, and then from random choice modality, as shown in \REF{bkm:Ref63461220}.

In order to answer the question as to how to account for the evaluative reading under transitive verbs, I assume that the evaluative reading ‘nonsense’ is a pragmatic implicature on \textit{cualquiera} or \textit{cualquier cosa} under transitive verbs that has been lexicalised under certain verbs such as \textit{mandar cualquiera} ‘talk nonsense’ in Argentinian Spanish. From the sentence in \REF{bkm:Ref63461360}, the hearer can infer that what the speaker said can be of the form p and not(p) at the same time. To utter a contradiction can be evaluated as nonsense. This is the most direct way to account for the meaning ‘nonsense’ of \textit{cualquiera}. Another, indirect, way to account for the meaning ‘nonsense’ is to derive ‘nonsense’ from the random choice meaning of \textit{cualquiera}; that is ‘something random’ can be interpreted as ‘nonsense’ in some pragmatic context. Namely when the agent who said something random did not reflect on what she/he said according to some judge. The result of the lack of reflection can be interpreted as something meaningless or as nonsense according the judge, such as in a situation when the agent was half-asleep. Thus, in this context the evaluator/judge can infer that the agent said something meaningless. Under this explanation, \textit{to talk} \textit{nonsense} would be derived from a contextual interpretation or pragmatic strengthening of the meaning of \textit{decir/mandar cualquiera} ‘to say something random’.

Under both analyses, the evaluative meaning of \textit{cualquiera} as ‘nonsense’ is a result of a pragmatic reasoning triggered by the use of FCIs under verbs of saying (see \citealt{Aloni2007} for indefinites in general, \citealt{Kellert2020semantic} on pragmatic reasoning in Mexican Spanish).

I find contexts that support my pragmatic analysis. The following example expresses that what the speaker actually answered does not correspond to what the speaker originally wanted to answer. In other words, the speaker’s intentions do not correspond to his/her actions/speech events. This is why \textit{cualquiera} has the meaning of ‘something unreflected’ in the following example:

\ea \label{ex:key:234} %ex 234
\gll No supe qué decirle, así que le mandé \textit{cualquiera}.\\
    not knew what say.her so that her sent \textsc{cualquiera}\\
    \glt ‘I did not know what to tell her, so I told her nonsense.’\\
    (@ 134)
\z

Another example, where the contextual interpretation restricts the interpretation of \textit{cualquiera} is given in the following sentence. Here the speaker has done everything except what she aimed to do, namely to fix the heating system. The exception is expressed through the contrast \textit{pero} ‘but’:

\ea \label{ex:key:235} %ex 235
    \gll Intenté  arreglar  la  calefacción,  pero  hice  \textit{cualquiera}.\\
    tried  fixing  the  heating.system  but  did  \textsc{cualquiera}\\
    \glt ‘I wanted to fix the heating system, but I did something out of place/nonsense (as a result of my attempt to fix the heating system).’\\
    (@ 135)
\z

In the following example \textit{decir} \textit{cualquiera} is interpreted negatively, similar to \textit{mandar la cagada} ‘to say bullshit’. It is interpreted as something unreflected or unreasonable.

\ea \label{ex:key:236} %ex 236
    \gll Che, ya te mandaste la cagada, dijiste \textit{cualquiera}, tiraste abajo a la mitad de los hombres diciendo que tienen que ser capaces de embarazar una mujer.\\
    that ya you said the bullshit said \textsc{cualquiera} pulled down to the half of the men saying that must that be capable of impregnate a woman\\
    \glt ‘Hey, you said bullshit, you said nonsense, you are demeaning half of the men by saying that they have to be capable of making a woman pregnant.’\\
    (@ 136)
\z

The inclusion of morally questionable alternatives like saying morally questionable things into the quantification domain of \textit{cualquiera} can result in a negative interpretation of \textit{decir} \textit{cualquiera} as saying all possible things including morally bad things:

\ea 
\label{ex:key:237} %ex 237
     \textit{decir cualquiera} \\
     ‘say anything possible including morally bad things’
\z

This negative meaning of \textit{cualquiera} is thus tantamount to FCI \textit{cualquiera} plus the pragmatic characterisation of the quantificational domain. The negative meaning of \textit{cualquiera} is the result of excluding certain alternatives from the domain of quantification which some evaluator considers as being good alternatives or of characterising the domain of quantification as consisting of all possible alternatives including morally questionable alternatives. I will therefore look into the domain quantification of \textit{cualquiera} in more detail in the next subsection.

\subsection{{Formal} {analysis} {of} {\textit{cualquiera}} {as} {‘nonsense’}} \label{sec:4.5.3}

I assume that the domain of quantification of ArgSp \textit{cualquiera} can be formally represented as a restriction of the quantification by an \textit{except}{}-phrase as in \REF{bkm:Ref63461516}. Assuming that \textit{except}{}-phrases contribute to the truth conditions of the sentence by restricting the domain of quantification of \textit{cualquiera} to the domain of predicates besides the predicate expressed by the \textit{except}{}-phrase (\citealt{Fintel1994}),\footnote{{The assumption that} {\textit{except}}{{}-phrases contribute to the truth conditions of the sentence they modify follows from the Restrictiveness Condition (\citealt{Fintel1994}): If a sentence with an} {\textit{except}}{{}-phrase is true, the same sentence without the} {\textit{except}}{{}-phrase is false (without making the exception, the sentence would be false).}} I must assume that \textit{cualquiera} in  \REF{bkm:Ref63461516} triggers universal quantification, which contributes to the truth conditions of the sentence it appears in:

\ea \label{bkm:Ref63461516} %ex 238   
‘I answered \textit{cualquiera} except what I wanted’

every([[\textit{thing}]] – [[\textit{what I wanted to answer}]])(\{x: I answered x\})
\z

\textit{Cualquiera} must be interpreted universally at some level of interpretation, because existential quantification does not allow \textit{except}{}-phrases (see \citealt{Fintel1994}, \citealt{Arregui2006}, \citealt{Gajewski2008}):

\ea \label{ex:key:239} %ex 239
  \#  I answered something except what I wanted.
\z

\textit{Except}{}-phrases of FCIs are a good argument that the universal inference of FCIs must be derived from a wider scope pragmatic operation or on the sentence level and not locally (see \sectref{sec:2.2} on such an analysis, where the universal interpretation is represented on the propositional level; see also \citealt{Chierchia2006}, \citealt{Gajewski2008}, but also \citealt{Dayal1998,Dayal2004}, who argues that FCIs are in fact universal quantifiers, with the result that \textit{except}{}-phrases operate locally). \textit{Except}{}-phrases should thus operate on a wide-scope universal implicature, that is, after the universal implicature of the FCI has been derived on the sentence or pragmatic level \citep{Gajewski2008}:\footnote{{The focus operator LEAST triggers exhaustivity of the alternative denoted by the exceptive; that is, it triggers the interpretation that the quantification domain of \textit{anything} includes all alternatives other than the one alternative subtracted from the domain of quantification.}}

\ea \label{ex:key:240} %ex 240
    Bill may have any flavour of ice cream but coffee.

    LF: [LEAST coffee]\textsubscript{X} Op\textsubscript{FC} [may [Bill have any flavour but X]]

    (\citealt{Gajewski2008}, \citealt{Gajewski2013}: 193)
\z

I assume that in case we have a pragmatically inferred \textit{except}{}-phrase as in \REF{bkm:Ref63461661}, repeated from \REF{ex:key:105}, \textit{cualquiera} is licensed by covert subtrigging (see \citealt{Quer2000} and \sectref{sec:2.1} for overt subtrigging):

\ea \label{bkm:Ref63461661} %ex 241
\ea \gll Respondí \textit{cualquiera}.\\
responded \textsc{cualquiera}\\
\glt ‘I answered anything that I could possibly answer except what I wanted to answer.’
\ex LF: [what I wanted to answer]\textsubscript{X} Op\textsubscript{FC} [I could possibly answer [I answered anything but X]]
\ex Prose: For every thing that I could possibly answer, I answered it, except what I wanted to answer.
\z
\z

This analysis explains the universal interpretation of \textit{cualquiera} on the one hand and the licensing of \textit{cualquiera} as FCI via subtrigging on the other.

Let me just briefly sketch an alternative analysis of \textit{cualquiera} in episodic contexts that relies on random choice modality and not on subtrigging. The analysis is very similar to the difference that the modal operator that licenses \textit{cualquiera} is some kind of modal, such as a counterfactual modal or a decision-based modal (see \sectref{sec:2.3}):

\ea
\label{bkm:Ref150787301} %ex 242
\ea \label{bkm:Ref150787301-1}
\gll Respondí \textit{cualquiera}.\\
responded \textsc{cualquiera}\\
\glt ‘I answered something random.’
\ex \label{bkm:Ref150787301-2} Existential component
‘I said something.’
\ex \label{bkm:Ref150787301-3} Counterfactual modal component: ‘And I could have said anything other than what I actually said.’

Decision-Based modal component: ‘And anything else that I did not actually say matches my decision to say something.’
\z
\z

The question now is how the evaluative characterisation of \textit{cualquiera} as ‘nonsense’ comes into play. Recall that this is the second challenge that I mentioned at the beginning of \sectref{sec:4.5.2}.

As already described in \sectref{sec:4.5.2}, I suggested that one could distinguish between contextual and lexicalised evaluative meaning. The contextual evaluative meaning of \textit{cualquiera} is better described as the evaluation of some judge in a given context. This means that \REF{bkm:Ref150787301-1} presupposes an attitude holder, the speaker of an utterance, who evaluates an event in which an agent has said something, though the agent could have said something different instead of what they actually chose to say (henceforth: event\textsubscript{said(agent)(anything possible)}). The attitude holder believes that this event is less preferred than a different event e' in which the agent said something else (henceforth: event\textsubscript{said(agent)(not(anything possible))}). This belief state of the attitude holder introduces an ordering scale of worlds that can be informally described as follows. All the worlds in which an agent said anything possible are less preferred than all worlds in which the agent did not say anything possible but said something specific/distinguished. I represent this ordering scale in line with \citet{Stalnaker1984} and \citegen{Heim1982} doxastic modal \textit{want}, which orders worlds according to the attitude holder’s preferences. \citet{Stalnaker1984} and \citet{Heim1982} suggest that \textit{want} is treated as a universal quantifier over worlds. ‘X wants S’ is interpreted as ‘X believes that a world w in which S is true is more desirable than a world in which S is false.’ Thus, \textit{want} is interpreted with respect to a doxastic modal base:

\ea \label{ex:key:243} %ex 243
    ⟦want⟧\textsuperscript{g,w} = λp ${\in}$ D\textsubscript{<s,t>} . λx ${\in}$ D\textsubscript{e}. Dox x,w $\subseteq$ \{w′: any world w‴ maximally similar to w′ such that w″ ${\in}$ p is more desirable to x in w than any world w‴ maximally similar to w′ such that w‴ ${\notin}$ p\}

    \citep[193]{Heim1982}
\z

I follow Heim’s analysis of attitude predicates such as \textit{want} and apply it to \textit{cualquiera}. On this analysis \textit{cualquiera} takes a proposition \textit{p}, an attitude holder x, and a doxastic modal corresponding to a set of worlds that are ordered according to x’s preferences in the actual world w in such a way as to represent p’s worlds being less preferred than non-p’s worlds:

\ea \label{bkm:Ref150787526} %ex 244

⟦cualquiera⟧\textsuperscript{g,w} = λp ${\in}$ D \textsubscript{<s,t>}. λx ${\in}$ D\textsubscript{e}. Dox <x,w> $\subseteq$ \{w′: any world that represent the negation of p is preferred to all worlds that represents p, i.e. the proposition that is the argument of \textit{cualquiera}\}

    (my adaptation of \citegen[193]{Heim1982} doxastic modality to \textit{cualquiera})
\z

If we apply this denotation to \REF{bkm:Ref150787301}, the ordered worlds correspond to any world in which the agent said anything possible is less preferred in w according to the attitude holder than the worlds in which the agent did not say anything possible, that is, in which the agent said something distinguished or particular.

The evaluative meaning component of \textit{cualquiera} thus corresponds to the attitude holder’s preferences in the actual world.

This analysis also captures the observation that the attitude holder does not have to be the speaker as shown in the following example (see also \sectref{sec:4.3.2}), by the fact that the attitude holder is represented as a variable in \REF{bkm:Ref150787526}.

\ea \label{ex:key:245} %ex 245
    \gll haría lo mismo pero si hay gente que no le gusta porque piensa que es \textit{cualquiera} (yo no pienso eso)...\\
do.\textsc{cond.3sg} it same but if there.is people that \textsc{neg} to.them likes because thinks that be \textsc{cualquiera} (I \textsc{neg} think that)\\
    \glt ‘I would do the same but there are people that don’t like it because they think it is \textit{nonsense} (I don’t think that)...’

    (@ 110)
\z

I assume that the lexicalisation of \textit{cualquiera} under \textit{mandar} can be represented as type-shifting in diachrony. Free Choice indefinite \textit{cualquiera} ‘anything’ of type <{}<e,t>,t> shifts into a property of type <e,t>, which corresponds to a noun that is interpreted as a collective thing that comprises all things that one could possibly say including contradictions. When used as an argument of \textit{mandar}, this collective thing can be paraphrased as ‘nonsense’.


\ea \label{ex:key:246} %ex 246
{[D° <{}<e,t>,t>} {\textit{cualquiera}}{ ‘anything’] $\rightarrow$ type-shifting [N° <e,t> ‘a collective thing that comprises all things that one could possibly say including contradictions’}]
\z

The analysis of \textit{cualquiera} (EVAL2) is thus the same as the one suggested in \sectref{sec:3.5} for \textit{cualquiera} in European Spanish (EVAL1), with the only difference that \textit{cualquiera} under \textit{mandar} in Argentinian Spanish is a property of things that denote propositions and not a property of kinds of people or entities as in EVAL1.

\subsection{{Summary} {of} {my} {analysis} {and} {discussion}} \label{sec:4.5.4}

The origin of the transitive \textit{cualquiera} as an FCI is analysed in parallel with \textit{cualquier cosa}, which is interpreted in certain pragmatic contexts as ‘anything possible except...’. I can now suggest the following diachronic steps that have led to the development of the evaluative meaning of \textit{cualquiera} under transitive verbs:

\ea \label{ex:key:247} %ex 247
  Diachronic steps of transitive \textit{cualquiera}
\begin{itemize}
  \item[Step 1:] Use of \textit{cualquiera} as ‘any’ to refer to a domain mentioned in the previous context in European Spanish and Argentinian Spanish. A special case of this use is the use of \textit{cualquiera} to refer back to previously mentioned \textit{cualquier cosa}.

    \item[Step 2:]  Use of \textit{cualquiera} in the sense of \textit{cualquier cosa} without previous mentioning of \textit{cualquier cosa} in Argentinian Spanish.

  \item[Step 3:]  Use of \textit{cualquiera} as \textit{cualquier cosa} in episodic contexts:

[\textsubscript{VP} \textit{decir/mandar}[episodic] [\textsubscript{DP} \textit{cualquiera=cualquier cosa}]] [\textsubscript{Adjunct NP} \textit{except} something specific or something the agent wanted to say]]]

  \item[Step 4:]  Evaluation of the thing or things that have been said according to preferences of an attitude holder: i.e. to say something randomly implies to say something without thinking they are less preferred according to an attitude holder than to say something distinguished or reflected.

  \item[Step 5:]  Interpretation of \textit{cualquiera} as ‘nonsense’.

  \item[Step 6:]  Lexicalisation of the pragmatic evaluation as ‘nonsense’ under \textit{mandar.}

    \item[Step 7:]  Syntactic shift from indefinite quantifier into a noun.
\end{itemize}
\z


As the diachronic development of \textit{cualquiera} under transitive verbs shows, the step towards lexicalisation of the evaluative meaning is the pragmatic interpretation of \textit{cualquiera} as something unreflected or something that can include contradictions according to some evaluator or judge. In contrast to \citet{Rizzosalierno2013}, I do not assume that the transitive \textit{cualquiera} has its origin in an analogy to the element \textit{cualunque}, because \textit{cualquiera} does not share the same distribution and frequency with \textit{cualunque}, but with \textit{cualquier cosa}. There is no occurrence with \textit{cualunque} under transitive verbs. However, we do find \textit{cualquier cosa} under transitive verbs.

The reanalysis of \textit{cualquiera} as an adjective in predicative position was already discussed in \chapref{sec:3} for European Spanish. The difference from European Spanish is that Argentinian Spanish allows degree modification of \textit{cualquiera} and the possibility of using \textit{cualquiera} as a predicative adjective without nominal modification, as shown in Step 2 (\tabref{tab:key:diachronic}).

\begin{table}
\begin{tabularx}{\textwidth}{lQ}
\lsptoprule
%\hhline%%replace by cmidrule{~--}

%\hhline%%replace by cmidrule{~--}
Step 1: & Use of \textit{cualquiera} as ‘any’ in postnominal position with the evaluative meaning:

 \textit{Juan es un hombre cualquiera.} (EurSp and ArgSp)

‘Juan is an ordinary man.’ COM\\
%\hhline%%replace by cmidrule{~--}
\tablevspace
Step 2: & Lexicalisation of COM in Argentinian Spanish and reanalysis as an adjective in +/– postnominal position

\textit{Juan es (un hombre) muy cualquiera.} (ArgSp)

‘Juan is a very ordinary man/Juan is very ordinary.’ COM\\
%\hhline%%replace by cmidrule{~--}
\lspbottomrule
\end{tabularx}
\caption{{Diachronic} {steps} {of} predicative {cualquiera}}
\label{tab:key:diachronic}
\end{table}

\chapref{sec:5} will offer a detailed diachronic investigation of the evaluative meaning (EVAL1) in diachrony. The diachronic analysis of EVAL2 will be tested in \sectref{sec:4.6.3}.

\section{{Further} {evidence} {for} {the} {analysis}} \label{sec:4.6}

In this section, further evidence for the analysis will be presented on the basis of synchronic and diachronic data.

\subsection{{Evidence} {from} {other} {languages}} \label{sec:4.6.1}

We find a similar reinterpretation of the free choice indefinite \textit{n’importe quoi} in French. As shown in \REF{bkm:Ref62641804}, the indefinite \textit{n’importe quoi} has different interpretations. It can have a random choice interpretation, as in \REF{bkm:Ref62641804-1}, or an evaluative interpretation, as in \REF{bkm:Ref62641804-2} (see \citealt{Larrivee2007}, \citealt{Paillard1997}, \citealt{Pescarini2009}, among others):

\ea \label{bkm:Ref62641804} %ex 249
French
\ea \label{bkm:Ref62641804-1} \gll {J’ai} {dit} \textit{n’importe quoi}\\
{I have} said \textsc{indefinite}\\

{(qui m’est passé par la tête, et ça s’est trouvé à être la bonne réponse!)}
\glt ‘I said something random (what came to my mind (first), and it happened to be the right answer).’ \citep[4]{Larrivee2007}
\ex \label{bkm:Ref62641804-2} 
\gll {Tu} {dis} \textit{n’importe quoi.} {Tu} {sais} {pas} {ce que} {tu} {dis.}\\
you say anything you know \textsc{neg} what you say\\
\glt ‘You’re talking nonsense. You don’t know what you’re saying.’
\glt ‘You’re talking bullshit/nonsense.’
    (@ 137)
\ex \label{bkm:Ref62641804-3} 
\gll {Tu} {dis} {du} {grand} \textit{n’importe quoi.}\\
you say of.the big anything\\
\glt ‘You’re talking bullshit/nonsense.’\\
    (@ 138)
\z
\z

If \textit{n’importe quoi} is used with a definite article and adjectival modification as in \REF{bkm:Ref62641804-3}, it is no longer ambiguous. It is then used as a noun with the evaluative interpretation of ‘nonsense’. \textit{N’importe quoi} shows similarities to \textit{cualquiera} under transitive verbs, as it can also be used as a noun with an evaluative interpretation. The difference between \textit{cualquiera} and \textit{n’importe quoi} is that \textit{n’importe quoi} cannot appear as a gradual adjective in a predicative position, whereas \textit{cualquiera} can as shown in \sectref{sec:4.3.2} (\textit{me parece re cualquiera ‘}it seems a complete bullshit to me’):

\ea[*]{{ \label{ex:key:250} %ex 250
\gll C’est très n’importe quoi. \\
{ it.is} very anything\\}\jambox*{(French)}
\glt {‘This is very bullshit.’}}
\z 

In contrast to \textit{n’importe quoi,} French FCI \textit{quelconque} can be modified with degree adverbs (see \citealt{Vlachou2007}). It has the meaning of ‘ordinary’:

\ea{ \label{ex:key:251} %ex 251
\gll Ce livre est assez \textit{quelconque}. \\
this book is very ordinary\\}\jambox*{(French)}
\glt {‘This book is rather ordinary.’}

{(@ 139)}
\z

The analysis of \textit{cualquiera} as a degree adjective applies to \textit{quelconque} as well.

A similar pattern was observed in Mexican Spanish for the lexical item \textit{equis} (originally meaning ‘the letter x’), which can be used as an indefinite with a free choice interpretation ‘any’ under conditionals as in \REF{ex:key:252}, but can also function as a gradable adjective that expresses EVAL1, as shown in \REF{ex:key:253} (see \citealt{Kellert2020semantic}):

\ea{ \label{ex:key:252} %ex 252
\gll Si por \textit{equis} razón no vas a venir, avísame  \\
if for some reason \textsc{neg} go to come let.me.know\\}\jambox*{(MexSp)}
    \glt ‘If you cannot come for some reason, let me know.’

    \glt ‘If for whatever reason you cannot come, let me know.’

    \citep[137]{Kellert2020semantic}
\z 

\ea{  \label{ex:key:253}  %ex 253
\gll un hombre bastante / demasiado / muy \textit{equis}  \\
a man quite / too / very average\\}\jambox*{(MexSp)}
    \glt ‘a very ordinary man (with no special qualities)’

    \citep[138]{Kellert2020semantic}
\z

More research needs to be done in order to testify the change of indefinites into evaluative nouns or adjectives in other languages (see \citealt{Simik2008} for Czech). In the future, I will also compare my results with more recent studies on the Mayan language Chuj (spoken in Guatemala and Mexico) that shows interesting parallels with the contact induced meaning variation of \textit{komon} ‘common, ordinary’ (\citealt{Alonso-OvalleRoyer2021}).

Another topic that needs further investigation is the licensing of indefinites in episodic contexts without overt subtrigging such as ArgSp \textit{respondí cualquiera} ‘I answered anything possible’ in other languages. \citet{Dayal1998} shows that \textit{any} can be licensed under perfective verbs without overt subtrigging in some English varieties (see \citealt{Dayal1998} for English varieties):

\ea \label{ex:key:254} %ex 254
\ea  Mary confidently answered any objections.

\ex After the dinner, I threw away any leftovers.
    \citep[446]{Dayal1998}
\z
\z 

In contrast to English, which rarely allows FCI \textit{any} in episodic contexts, Romance indefinites such as \textit{cualquiera} do allow FCIs in episodic contexts (see \sectref{sec:2.3}). One possible explanation why \textit{cualquiera} appears more often in episodic contexts than English \textit{any} is that \textit{cualquiera} lexically encodes modality due to its relative clause origin from free choice relative \textit{qual quier} ‘whatever you want’ (see \citealt{Rivero2011} and \sectref{sec:2.2}). This explanation also holds for French as \textit{n’importe quoi} lexically encodes a modal meaning ‘it is not important what’. It seems that English \textit{any} does not encode modality lexically and thus does not allow \textit{any} to appear in episodic contexts as in Spanish and French.

\subsection{{Evidence} {from} {comparison} {with} {\textit{cualquier} \textit{cosa}}} \label{sec:4.6.2}

My analysis of ArgSp \textit{cualquiera} is based on the FCI \textit{cualquiera}. It is thus possible that \textit{cualquier cosa} has undergone a similar change. I will therefore compare the two items and their uses. All data is taken from CdE\_web dialects. As I will show \textit{cualquier cosa} is on its way to reanalysis as a degree adjective ‘bad/out of place’.

 \textit{Cualquier cosa} can be used predicatively on par with \textit{cualquiera} and have the same evaluative meaning of ‘bad, incorrect, out of place’. This is shown by the use of \textit{cualquier cosa} and \textit{cualquiera} under predicative verbs such as \textit{ser/parecer} ‘be’/‘behave’.

\ea \label{ex:key:255} %ex 255
    \gll Ese estudio es \textit{cualquier} \textit{cosa}, realmente \textit{cualquier} \textit{cosa}.\\
    that study is \textsc{cualquier} thing really \textsc{cualquier} thing\\
    \glt ‘This study is bad/out of place, really bad/out of place.’\\
    (@ 140)
\z

As with \textit{cualquiera}, \textit{parece cualquier cosa} is often accompanied by some evaluator (e.g. ‘I find’):

\ea \label{ex:key:256} %ex 256
    \gll Me \textit{parece} \textit{cualquier} \textit{cosa} la designacion de este día, francamente.\\
    me seems \textsc{cualquier} thing the appointment of this day frankly\\
    \glt ‘Today’s appointment seems out of place to me, frankly.’\\
    (@ 141)
\z

\ea \label{ex:key:257} %ex 257
    \gll Me \textit{parece} \textit{cualquier} \textit{cosa} lo que planean hacer con la línea H. No hay estudios que demuestren la factibilida del neuvo  recorrido.\\
    to.me seems \textsc{cualquier} thing the that plan to.do with the line H \textsc{neg} have studies that show the feasibility of.the new route\\
    \glt ‘I find it nonsense what they are planning to do with the H line.’\\
    (@ 142)
\z

A new finding is that \textit{cualquier cosa} can even be used as a degree adjective (only one occurrence has been identfieid in CdE\_web dialects with degree adverb \textit{muy}):

\ea \label{ex:key:258} %ex 258
    \gll No la soporto en ninguna. Verónica Coincido con Maid en esa pareja es \textit{muy} \textit{cualquier} \textit{cosa!}\\
    \textsc{neg} her tolerate in no-one Verónica Coincido with Maid in that couple is very \textsc{cualquier} thing\\
    \glt ‘I can’t stand any of them. Verónica Coincido with Maid in that couple is completely out of place.’\\
    (@ 143)
\z

Web research has turned up more examples with degree modification by \textit{re/muy} from young bloggers or Facebook users:

\ea \label{ex:key:259} %ex 259
    \gll Mi pose \textit{es} \textit{re} \textit{cualquier} cosa pero la edit es preciosa!\\
    my pose is very \textsc{cualquier} thing but the edit is gorgeous\\
    \glt ‘My pose is completely out of place, but the edit is gorgeous!’\\
    (@ 144)
\z

\ea \label{ex:key:260} %ex 260
    \gll Ya es \textit{muy} \textit{cualquier} cosa ya.\\
    now is very \textsc{cualquier} thing now\\
    \glt ‘It is completely out of place, it definitely is.’\\
    (@ 145)
\z

\ea \label{ex:key:261} %ex 261
    \gll Julio Ponce suena \textit{muy} \textit{cualquier} cosa,\\
    Julio Ponce plays very \textsc{cualquier} thing\\
    \glt ‘Julio Ponce plays completely out of place.’\\
    (@ 146)
\z

I also found \textit{medio} with \textit{cualquier cosa} by young forum users and bloggers:

\ea \label{ex:key:262} %ex 262
    \gll Creo que te podés presentar de la nada, pero conviene mandarle un mail. Es \textit{medio} \textit{cualquier} \textit{cosa} {el} sistema (y mejor ni hablar  del adelanto de la segunda fecha todavía estoy con la vena).\\
    believe that you can appear from the nothing but is.advisable send.him a email is half \textsc{cualquier} thing the system (and better not.even talk of.the forwarding of the second date still am with the mood)\\
    \glt ‘I think that you can go there directly, but it is advisable to send her/him an email. The system is somewhat out of place (and it’s better if I don’t even talk about moving the second date forward, I am still in a bad mood).’\\
    (@ 147)
\z

\ea \label{ex:key:263} %ex 263
    \gll Sí, es cierto, la actuación de Cecilia Dopazo es \textit{medio} \textit{cualquier} \textit{cosa} […] y fue pésima la elección de ella para el papel.\\
    yes is true the performance of Cecilia Dopazo is half \textsc{cualquier} thing {} and was awful the decision of her for the role\\
    \glt ‘Yes, it’s true, Cecilia Dopazo’s performance was somewhat out of place […] and choosing her for the role was an awful decision.’\\
    (@ 148)
\z

These data show that \textit{cualquiera} and \textit{cualquier cosa} can function as gradable adjectives and be modified by different degree modifiers (\textit{muy, re, medio}) that denote +/– high degrees for some speakers. However, more research is needed to identify the sociolinguistic parameters in the acceptability of degree modification of \textit{cualquier cosa}.

With respect to degree suffixes, I also found a difference in interpretation between \textit{cualquiera} and \textit{cualquier cosa}. The diminuitive suffix \textit{-ita} can combine with \textit{cualquiera}, and in this case the speaker attenuates the bad quality ‘a bit bad’ (see \sectref{sec:2.1}):

\ea \label{ex:key:264} %ex 264
    \gll Bueno, cerramos este post que {es} \textit{cualquierita}.\\
    well close this post that is \textsc{cualquier}.\textsc{dim}\\
    \glt ‘Well, let’s close this post because it’s a bit nonsensical.’\\
    (@ 118)
\z

The noun \textit{cosa} can be modified by a diminuitive -\textit{ita} to denote a little tiny thing (a \textit{minimiser}, see \citealt{Krifka1995}), and the combination \textit{cualquier cosita} denotes an FCI ‘any little tiny thing’:

\ea \label{ex:key:265} %ex 265
    \gll Dando educación de calidad y no \textit{cualquier} \textit{cosita}.\\
    offering education of quality and \textsc{neg} \textsc{cualquier} thing.\textsc{dim}\\
    \glt ‘Offering quality education and not just any little thing.’\\
    (@ 149)
\z

\ea \label{ex:key:266} %ex 266
    \gll de \textit{cualquier} \textit{cosita} llora peor que un chico\\
    of \textsc{cualquier} thing.\textsc{dim} weeps worse that a boy\\
    \glt ‘He weeps over any little thing, worse than a little boy.’\\
    (@ 150)
\z

Whereas \textit{cualquier} in \textit{cualquier cosa} is originally a quantificational determiner D° with an FC meaning ‘anything’ (both in European and in Latin American Spanish), as shown in \REF{bkm:Ref150787837}, in Argentinian Spanish the whole string \textit{cualquier cosa} has additionally developed a special evaluative meaning and has also been recategoried as an Adj°, see \REF{bkm:Ref150787813}:

\ea \label{bkm:Ref150787837} %ex 267
\ea 
D° cualquier cosa = ‘anything’ 
\ex λP.Ǝx. thing(\textit{x}) \& P(x) + Free Choice Implicature
(ArgSp and EurSp)
\z 
\z

\ea \label{bkm:Ref150787813}   %ex 268 
Adj° cualquier cosa = λd.λx. bad(d)(x)

(ArgSp)
\z


The observation that \textit{cualquiera} but not \textit{cualquier cosa} has been found with comparatives and the observation that \textit{cualquier cosa} is less frequent with gradable adjectives seems to suggest that the process of recategorisation into a gradable adjective is more advanced with \textit{cualquiera} than with \textit{cualquier cosa}. If the difference between \textit{cualquiera} and \textit{cualquier cosa} lies in the difference in the ongoing process of lexicalisation, we should observe differences in age distribution or geographical distribution. I will test this hypothesis with native speaker judgements in future research.

\subsection{{Diachronic} {evidence}} \label{sec:4.6.3}

For diachronic research of \textit{cualquiera} and \textit{cualquier cosa} in Argentinian Spanish, I first used the diachronic corpus CORDIAM (Corpus Diacrónico y Diatópico del Español, CORDIAM, \href{http://www.cordiam.org}{www.cordiam.org}, which contains a total of 10,889,176 words from 19 Latin-American countries from 1494 to 1905, with documents from archives, journals, and literature. It is annotated for geographical origins. The corpus CORDIAM contains only 104 instances of \textit{cualquiera} in Argentinian Spanish, only one instance of N \textit{cualquiera}, and zero instances of \textit{cualquier cosa} and \textit{cualquiera} with the meaning of \textit{cualquier cosa}. The remaining instances use \textit{cualquier} NP and \textit{cualquiera} as ‘any NP’ or ‘anyone’:

\ea \label{ex:key:269} %ex 269
    \gll Ella [...], á este le dirije un cumplimiento, á ese una mirada simpática, á aquel un apóstrofe \textit{cualquiera},\\
she [...] to this.one to.him directs a greeting to that.one a look friendly to that.one a apostrophe \textsc{cualquiera}\\
    \glt ‘She sends this one a compliment, to the other one a nice smile, and to that one a random insult.’

{    (19}{{th}}{, anonymous, journal:} {\textit{La Camelia}}{; ArgSp)}
\z

I also used a larger diachronic corpus, CORDE (Corpus Diacrónico del Español), which is a diachronic corpus with written texts spanning from Old Spanish to 1974. It contains 250 million written words from different genres, types, and geographical origins. \textit{Cualquiera} appears 535 times in Argentinian Spanish starting from the 18th century. Before that, the form \textit{qualquiera} appeared nine times from the 16th to the 18th century (see \chapref{sec:5} on different forms). There is no case of gradable \textit{cualquiera} in predicative position (e.g. \textit{es muy cualquiera}). However, there is one occurrence of predicative \textit{cualquiera} without nominal modification with the meaning ‘nothing important’ or ‘ordinary’:

\ea \label{bkm:Ref150787886} %ex 270
\textit{Cualquiera} as a predicative complement

    \gll {Su} {ocupación} {era} \textit{cualquiera.}\\
    his position was \textsc{cualquiera}\\
    \glt ‘He had an ordinary position./His position was nothing important.’\\
{(1926, Güiraldes,} {\textit{Don Segundo Sombra}}{; ArgSp)}
\z

One interesting observation is that \textsc{un} N \textit{cualquiera}, such as in \REF{bkm:Ref150787886}, appears more often in affirmative contexts in Argentinian Spanish than it does in European Spanish (56\% in Argentinian Spanish vs 18\% in EurSp). This observation of the relatively high frequency of \textsc{un} N \textit{cualquiera} in Argentinian Spanish can be interpreted as an indicator that \textit{cualquiera} is undergoing the process of lexicalisation of the meaning ‘ordinary’ and ‘bad’. In \chapref{sec:5}, we will look at the development of \textit{cualquiera} in European Spanish and investigate the appearance of the evaluative reading (EVAL1).

There is no diachronic data that could confirm Step 2 in \sectref{sec:4.5.3}. CORDE does not contain \textit{cualquiera} as a substitute for \textit{cualquier cosa}. CORDE contains 167 examples of \textit{cualquier cosa} in Argentinian Spanish among which five examples of \textit{cualquier cosa} are embedded in episodic contexts under verbs of saying (\textit{decir}, \textit{hablar de, murmurar, conversar de}). Note that \textit{cualquier cosa} does not mean ‘nonsense’ in these contexts, but rather has a weaker reading of ‘nothing in particular’ or ‘something unimportant’; that is, the speaker does not care to specify the reference of the object of conversation, as in \REF{ex:key:271-1}, or the speaker has no plan in mind of what he was going to do, as in \REF{ex:key:271-2} (see the indifference or indiscriminative reading in 2.4).

\ea \label{ex:key:271} %ex 271

  \textit{Cualquier cosa} in Argentinian Spanish
\ea \label{ex:key:271-1} \gll Y el público, formado por la gente de huella y de estancia, conversaba de \textit{cualquier} \textit{cosa}.\\
and the audience, formed by the people of footprint and of stay, conversed about \textsc{cualquier} thing\\
    \glt ‘The general public, herders and ranchers, hummed with talk about everything and nothing.’

    (1926, Güiraldes, \textit{Don Segundo Sombra}; ArgSp)
\ex \label{ex:key:271-2} \gll Yo salí muy decidido a \textit{cualquier} \textit{cosa}, pero a nada en particular.\\
I went.out very determined to \textsc{cualquier} thing, but to nothing in particular\\
    \glt ‘I was not sure what I was going to do.’

    (1940, Bioy Casares, \textit{La invención de Morel}; ArgSp)
\ex \label{ex:key:271-3} \gll Y de pronto, para romper el intolerable silencio, yo decía \textit{cualquier} cosa que en otro tiempo podía haber tenido interés para el tipógrafo. - La fora ha declarado una huelga de estibadores.\\
and of suddenly to break the intolerable silence I said \textsc{cualquier} thing that in another time could have had interest for the typographer. - the union has declared a strike of stevedores\\
    \glt ‘And suddenly, in order to break the unbearable silence, I said just something or other that in another time would have had some interest for the typographer.’

    (1961, Sábato, \textit{Héroes y tumbas}; ArgSp)
    
\ex \label{ex:key:271-4}
\gll Estábamos tendidos de espaldas, uno al lado del otro, y ella me acariciaba maquinalmente, yo tenía frío y ella me hablaba de \textit{cualquier} \textit{cosa}, de la pelea que acababa de ocurrir en el bar, de las tormentas de marzo, […]\\
were lying on backs, one to.the side of.the other and she me caressed mechanically, I had cold and she me spoke about \textsc{cualquier} thing, about the fight that had.just of occur in the bar, of the storms of March\\
    \glt ‘We were close together, one against the other, and she caressed automatically, I was cold and she talked to me about something or other, about the fight that had just taken place in the bar, about the storms in March, […] .’

    (1963, Cortazár, \textit{Rayuela}; ArgSp)
\ex \label{ex:key:271-5} \gll Suspirando, Traveler se volvió a la cama. A una pregunta soñolienta de Talita, le acarició el pelo y murmuró \textit{cualquier} \textit{cosa}. Talita besó el aire, manoteó un poco, se tranquilizó.\\
sighing Traveler \textsc{refl} turned to the bed to a question sleepy of Talita to.her caressed the hair and murmured \textsc{cualquier} thing. Talita kissed the air, waved a little, \textsc{refl} calmed.down\\

    \glt ‘Sighing, Traveler went back to his room. After Talita asked him a question in a sleepy state, he caressed her hair and whispered something. […]’

    (1963, Cortazár, \textit{Rayuela}; ArgSp)
\z
\z

{The ratio of meanings per century is summarised in \tabref{tab:4:extra}}.

\begin{table}
\begin{tabularx}{\textwidth}{ lrQ@{} }
\lsptoprule
 Linguistic context & CORDIAM & CORDE \\ 
 \midrule
 bare \textit{cualquiera} & 0 & 1 (EVAL) \\
 \textit{cualquier cosa} & 0 & 167 (162 with FCI, 5 with EVAL under verbs of saying in 20th c.) \\
 other context & 103 (FCI) & 367 (FCI) \\
 \midrule
 total & 103 & 535\\
 \lspbottomrule
\end{tabularx}
\caption{\label{tab:4:extra} Diachronic distribution of \textit{cualquiera} and \textit{cualquier cosa} in CORDIAM and CORDE}
\end{table}

To sum up, the diachronic corpora contain only one example of a predicative adjectival use of \textit{cualquiera} and no adjectival use of \textit{cualquier cosa}. They do not show the special ‘nonsense’ reading under verbs of saying. The first examples that appear with degree modification are examples from the modern use of Argentinian Spanish in CdE\_web dialects in the 21st century, as reported in \sectref{sec:4.3}. The earliest occurrences that I could find of \textit{decir/mandar cualquiera} originate from the 21st century, although according to \citet{Conde2011} \textit{decir/mandar cualquiera} is apparently from Lunfardo, which must be older than the 21st century. One possible explanation why there is no diachronic data of special uses of \textit{cualquiera} in Argentinian Spanish is that diachronic data is mostly taken from written language, whereas the special uses of \textit{cualquiera} represent spoken language. Another explanation is that the use of \textit{cualquiera} as a gradable adjective with the meaning ‘ordinary’ and ‘bad’ is a recent innovation of native speakers of Argentinian Spanish. Describing diachronic change of spoken language is thus extremely difficult with corpus methods due to the scarcity of corpus data. 

{In the future, single texts representing a higher degree of communicative immediacy, ``Nähesprache'' in the sense of \citet{koch1990gesprochene}, such as letters or personal journals and drama will be considered.}

